\documentclass[11pt]{article}
\usepackage[T1]{fontenc}
\usepackage[utf8]{inputenc}
\usepackage{lmodern,microtype,geometry,amsmath,amssymb,amsthm,mathtools}
\usepackage[colorlinks=true,linkcolor=blue,citecolor=blue,urlcolor=blue]{hyperref}
\geometry{margin=1in}
\setlength{\emergencystretch}{3em}
\newtheorem{theorem}{Theorem}
\newtheorem{lemma}[theorem]{Lemma}
\newtheorem{corollary}[theorem]{Corollary}
\newtheorem{remark}[theorem]{Remark}
\newcommand{\E}{\mathbb E}
\newcommand{\R}{\mathbb R}
\newcommand{\Prb}{\mathbb P}
\newcommand{\norm}[1]{\left\|#1\right\|}
\title{Quartic-width learning bounds for robust Bayesian tanh networks}
\author{ChatGPT-6 Astra \and Claude Fable 5 and Opus 5 \and Kaarel H\"anni \and Dmitry Vaintrob\thanks{Authors are listed in alphabetical order.}}
\date{Research draft, 29 October 2026}
\hypersetup{pdftitle={Quartic-width learning bounds for robust Bayesian tanh networks},pdfauthor={ChatGPT-6 Astra; Claude Fable 5 and Opus 5; Kaarel Hanni; Dmitry Vaintrob}}
\begin{document}
\maketitle
\begin{abstract}
We prove sample bounds for Gaussian Bayesian tanh networks restricted to
predictions stable under simultaneous internal disturbances. For a layered
Boolean teacher of width $C$ and $L-1$ gate stages, a sufficient sample
count is $\widetilde O(C^4L^9)$, uniformly over every learner width
$N\ge C$ and every mean-field-or-stronger scale $\gamma\ge1$.
For $\gamma\ge7/4$, the bound improves to $\widetilde O(C^4L^8)$.
At $\gamma\ge2$, and in the strong-prior limit at each fixed width,
a separate argument gives
$\widetilde O(\min\{C^4L^7,C^{9/2}L^6\})$. The same stronger bound
also holds at a positive scale ceiling independent of learner width.
The robustness restriction is unchanged throughout. Input dimension,
robustness calibration, accuracy, and confidence are fixed in these
rate displays. The all-scale proof combines entropy-weighted regression
with a Gaussian-weighted description of uniformly sampled coefficient
rows. A spectral truncation separates a bounded-energy network from a
background of small operator norm and yields the uniform cold-scale
bound. Multilevel sampling supplies the depth-six tradeoff.
We give exact budgets and proofs, and explicitly state the previously
established reconstruction estimates used. These results are conventional
mathematical proofs, not Lean formalizations.
\end{abstract}

\section{Introduction}

A polynomial sample bound can still be very large. For a neural network
learning a Boolean circuit, it matters how the bound depends jointly on
the circuit's width and depth. We study exact Gaussian Bayesian prediction
with tanh hidden layers and a linear readout, conditioned on noiseless
training margins. The prior retains only networks whose clipped outputs
are stable under a calibrated ball of simultaneous internal disturbances.
The restriction is stated in full below and is fixed before observing
training labels.

The main result is $\widetilde O(C^4L^9)$ samples throughout the
mean-field-or-stronger Gaussian regime, at every learner width $N\ge C$.
For $\gamma\ge7/4$, a separate interpolation gives
$\widetilde O(C^4L^8)$. At $\gamma\ge2$, the bounds are
$\widetilde O(C^4L^7)$ and $\widetilde O(C^{9/2}L^6)$; one may choose
the smaller. The latter improves depth dependence at a cost in the
width exponent. All statements allow an arbitrary common iid training
and test input distribution.

The all-scale construction measures approximation error using a matrix
that records conditional Gaussian entropy in different active
directions. Independent-batch regression controls error in that metric.
An upper estimate of the total variance permits a stable nonlinear
projection and one Gaussian coefficient per reconstructed row. The
resulting comparison functions are finite arithmetic programs; they
need not be neural networks. Only the original posterior networks must
satisfy the robustness restriction.

A substantial saving comes from the probabilities assigned to these
finite descriptions. Uniformly sampled rows inherit only their fraction
of a full Gaussian block's relative entropy. We use their quantized
Gaussian probabilities in the comparison prior, and encode the remaining
parameters separately. This is a proof device: the actual Bayesian
learner's prior is unchanged. Direct posterior bounds handle widths
where constructing the code would cost too much, and a finite mixture
of reference priors gives one event valid for all widths and scales.

The stronger bounds use a different idea. Clamping each internal
matrix's singular values at a constant times the Gaussian scale leaves
a residual network whose energy has a width-independent posterior tail.
The clamped part has small operator norm. When the scale is sufficiently
small, discarding it uses only part of the allowed internal disturbance.
This gives a positive cold-scale threshold independent of learner width,
and interpolating it with a direct Bayesian bound gives the intermediate
$\gamma\ge7/4$ result. Multilevel sampling of two consecutive activation
layers gives the separate depth-six tradeoff. No merging, coherence,
or spectral-product condition is added to the prior.

After the model and main statements, we give the simpler energy-based
code and its posterior argument. The appendices give the target
cylinder and supporting reconstruction estimates, followed by
entropy-metric compression (Appendix~\ref{app:entropy-metric}),
Gaussian-weighted coding and its all-scale learning theorem
(Appendices~\ref{app:weighted-entropy}--\ref{app:weighted-code}),
multilevel compression (Appendix~\ref{app:mlmc-code}), and the
spectral cold-scale bound and interpolation
(Appendices~\ref{app:spectral-cold}--\ref{app:spectral-interpolation}).
The supporting sequential and factored codes build on~\cite{sequential}.
No efficient posterior inference algorithm or optimality of the sample
powers is asserted.

\section{Learner, target, and robustness restriction}

Let $d\ge1$, $L\ge2$, and $H=L-1$. Inputs lie in
$\mathcal X=\{-1,1\}^d$, with arbitrary distribution $D$.
A learner of hidden width $N$ has parameters $w=(W^1,\ldots,W^H,v)$ and
computes
\[
 h^1=\tanh(W^1\bar x),\qquad
 h^s=\tanh(W^sh^{s-1}),\qquad
 z_w(x)=v^\top h^H,\qquad \bar x=(x,1).
\]
The raw law $\mu_{N,t}$ has independent centered Gaussian entries of
variance $t^2/d$ in $W^1$, $t^2/N$ in the internal matrices, and
$t^2/N^2$ in $v$. Here $0<t\le1$. The scale may be any prescribed
function $t=t(d,L,N)$; it need not be constant as width varies.
In particular $t(N)=N^{(1-\gamma)/2}$ gives variances
$N^{1-\gamma}/d$, $N^{-\gamma}$, and $N^{-\gamma-1}$.

Define weighted energy, readout energy, and clipped score by
\begin{align*}
 \kappa_N(w)&=\frac dN\norm{W^1}_F^2+
   \sum_{s=2}^H\norm{W^s}_F^2+N\norm v_2^2,\\
 E_v(w)&=N\norm v_2^2,\qquad
 T(u)=\max\{-1/2,\min\{1/2,u\}\}.
\end{align*}
Disturbances are added after each activation and at the readout:
\[
 h^1_\eta=\tanh(W^1\bar x)+\eta^1,\quad
 h^s_\eta=\tanh(W^sh^{s-1}_\eta)+\eta^s,\quad
 z_w(x;\eta)=v^\top h^H_\eta+\eta^{\rm out}.
\]
Their norm is
\[
 \norm\eta_*^2=\sum_{s=1}^H\frac{\norm{\eta^s}_2^2}{N}
                  +|\eta^{\rm out}|^2.
\]
Fix $c>0$ before seeing labels. The support condition is
\begin{equation}\label{eq:support}
 \mathcal R_{N,c}=\left\{w:
 |T(z_w(x;\eta))-T(z_w(x))|\le\frac1{16}
 \text{ for all }x\in\mathcal X,\quad
 \norm\eta_*\le\frac c{\sqrt{1+HE_v(w)}}\right\}.
\end{equation}
Noisy hidden values are allowed outside $[-1,1]^N$.
This is a label-independent Borel condition, since the input set is finite
and the response is continuous in weights and disturbances.

For a sample $S=((X_i,f(X_i)))_{i=1}^n$ with iid $X_i\sim D$, let
$Q_{S,N,t,c}$ be the raw Gaussian conditioned on
\begin{equation}\label{eq:conditioning}
 w\in\mathcal R_{N,c},\qquad
 f(X_i)z_w(X_i)>1/2\quad(1\le i\le n).
\end{equation}
There is no initial clean-output-margin condition on unobserved inputs.
The risk is posterior-average test-margin failure,
\[
 R(Q)=\E_{w\sim Q}\Prb_{X\sim D}\{f(X)z_w(X)\le1/4\}.
\]

The target is a single-output Boolean circuit of width $C\ge2$ and $H$
layered stages, using NAND, COPY, and NOT. A constant-one wire is reserved
at every stage and counted toward $C$; compilation overhead from another
basis is also counted. Unused slots and stages may contain constants or
copies. Thus a teacher with $h\le H$ stages can be padded to the learner's
depth. If the learner matches the teacher's depth, $L=h+1$; if it is deeper,
its actual $L$ occurs in the bound. Width independence does not mean depth
independence.

\section{Learning results}

\begin{theorem}[Quartic width throughout the good Gaussian regimes]
\label{thm:headline-all-scales}
Fix $d,c,\varepsilon,\delta$ with $c>0$ and $0<\varepsilon,\delta<1$.
Put $F_0=1+\log^2(2+C)$ and $G=\log(2+CL)$.
There is a constant $A_0$, depending only on these four quantities,
such that
\[
 n\ge A_0 C^4L^9F_0^4G^2
\]
implies, with sample probability at least $1-\delta$,
$R(Q_{S,N,t,c})\le\varepsilon$ simultaneously for every $N\ge C$
and every $0<t\le1$. In particular, this includes every
$\gamma\ge1$.

The statement holds for every Boolean teacher and iid input law
specified above. It also holds after adding the label-independent
prior restriction $|z_w(x)|>1/2$ at every Boolean input.
\end{theorem}

\begin{proof}[Proof and exact-budget specification]
Use the product target cost in \eqref{eq:product-cylinder-cost} and
the posterior budgets in \eqref{eq:factor-posterior-budgets}.
The readout-tail calculation in Appendix~\ref{app:factor-code} retains
posterior mass at least $1-\beta$ on a common robustness radius $r$.
Its retained law $Q'$ satisfies $Q'\le \exp(NK/t^2)\mu_{N,t}$.
Theorem~\ref{thm:weighted-code-kl} supplies a comparison law whose
uniformly sampled coefficient rows are charged by their Gaussian
relative entropy. The remaining parameters have a finite description
of order $LK^2/e^2$, apart from logarithms.
The scale and width cases, their common statistical event, and their
exact budgets are in Theorem~\ref{thm:weighted-all-scales}.
Corollary~\ref{cor:weighted-powers} substitutes
\[
 K=O(CLF_0),\qquad r^{-2}=O(CL^2F_0),\qquad
 e^{-1}=O(CL^3F_0),\qquad \eta^{-1}=O(L)
\]
and proves the displayed logarithmic factor as well as the
simultaneous quantifiers. The same target cylinder survives optional
global-margin conditioning, so the denominator, tail, and comparison
arguments apply to that variant too.
\end{proof}

\begin{corollary}[Intermediate Gaussian exponents]
\label{cor:headline-intermediate}
At fixed $d,c,\varepsilon,\delta$, a sufficient sample count is
\[
 O\!\left(C^4L^8F_0^4\log^2(2+L)G\right)
\]
simultaneously for every $N\ge C$ and every
$0<t^2\le N^{-3/4}$, hence for every $\gamma\ge7/4$.
More precisely, for each fixed $\gamma>1$ the sufficient order is
\[
 \widetilde O\!\left(\min\left\{
 C^4L^9,\quad
 C^4L^7+C^{2+1/(\gamma-1)}L^{4+3/(\gamma-1)}
 \right\}\right).
\]
The prediction convention and optional margin restriction are those
of Theorem~\ref{thm:headline-all-scales}.
\end{corollary}

The proof, exact statistical budget, and interpretation of this
interpolation are in Corollary~\ref{cor:spectral-interpolation}.
At $\gamma=7/4$, the two-term order before simplifying the width
power is $\widetilde O(C^4L^7+C^{10/3}L^8)$.

\begin{theorem}[A further depth improvement at stronger scales]
\label{thm:headline-mlmc}
At fixed $d,c,\varepsilon,\delta$, the same robust Bayesian learner
has sufficient sample order
\[
 \widetilde O\!\left(\min\{C^4L^7,C^{9/2}L^6\}\right)
\]
simultaneously for every $N\ge C$ and $0<t\le N^{-1/2}$,
including every $\gamma\ge2$. The same order suffices in the
strong-prior limit at each fixed width. Both conclusions refer to
posterior-average test-margin failure $R(Q)\le\varepsilon$ with
sample probability at least $1-\delta$.

The same sample order holds simultaneously for all $N\ge C$ at every
positive scale $t^2\le u_c$, where
$u_c=r_{\rm sp}^2/[256(L-1)]$ and $r_{\rm sp}$ is defined in
\eqref{eq:spectral-budget}. This ceiling depends on teacher resources,
calibration, and accuracy, but not on learner width.
\end{theorem}

The finite exponent alternatives and proof are in
Theorem~\ref{thm:mlmc-posterior}; the positive width-independent
ceiling is proved in Theorem~\ref{thm:spectral-cold}. The depth-six branch improves the
powers when $C<L^2$. The baseline quartic code and its probability
calculation are given next; the multilevel construction replaces only
that code.

\subsection{The energy-based baseline and explicit stronger scales}

The following separate argument trades its stronger scale hypothesis for
quartic dependence on teacher width. Define
\begin{align}
 e_*&=\frac1{16\sqrt{2C}},\qquad \tau_0=e_*/4,\notag\\
 b_L&=\sqrt{16\log9+8\log(4L)},\notag\\
 A&=\max\{32b_L,2\log(8/e_*)\},\qquad a=\max\{2,256c\},\notag\\
 \rho&=\min\{1/4,A/[4(d+1)],Ae_*/(64\sqrt C)\},\notag\\
 q_0&=d+2+(L-2)C,\notag\\
 B_0&=9dA^2+18(L-2)CA^2+2a^2C,\notag\\
 D_*&=B_0+q_0\log\frac{\exp(1)\sqrt{2\pi}}{2\rho}+\log2.
 \label{eq:target-cost}
\end{align}
These constants provide a target-realizing subset of the prior support of
raw mass at least $\exp(-ND_*/t^2)$ at every $N\ge C$; the proof is in
Appendix~\ref{app:cylinder}. In particular $D_*\ge1$ and every posterior
in the theorem below is well-defined.

For $K\ge1$, set
\begin{align}
 r&=\min\{1/6,c/\sqrt{1+HK}\}, &
 m&=\left\lceil256Kr^{-2}[d+\log(400L)]\right\rceil,\notag\\
 B&=\sqrt{4mK}, &
 \zeta&=\frac r{4\sqrt{(d+1)^2+Hm^2}},\notag\\
 P&=m(d+1)+(L-2)m^2+m, &
 J&=P\log(3+2B/\zeta).
 \label{eq:budget}
\end{align}

\begin{theorem}[Uniform-width learning with polynomial coding precision]
\label{thm:main}
Let $0<\varepsilon,\delta<1$ and suppose $f$ has the Boolean implementation
specified above. The sample condition
\begin{equation}\label{eq:samples}
 n\ge\frac2\varepsilon\left[J+\log(1/\delta)\right]
\end{equation}
suffices in either of the following cases:
\begin{enumerate}
\item Set $K=8D_*$. Allow every $N\ge C$ and every positive scale satisfying
\[
 t^2\le\min\left\{1,
 \frac{ND_*}{(P_N/2)\log2+\log(2/\varepsilon)}\right\},
 \qquad P_N=N(d+1)+(L-2)N^2+N.
\]
\item Set $K=8D_*+4(d+L)+4\log(2/\varepsilon)$. Allow every $N\ge C$
and $0<t\le N^{-1/2}$, including every fixed exponent $\gamma\ge2$.
\end{enumerate}
With probability at least $1-\delta$ over the sample, the risk
$R(Q_{S,N,t,c})$ is at most $\varepsilon$ simultaneously for all indicated
widths and scales.
\end{theorem}

For every fixed width, the first clause permits all sufficiently small
positive scales and therefore gives the strong-prior risk conclusion
$\limsup_{t\downarrow0}R(Q_{S,N,t,c})\le\varepsilon$ on the same sample
event. No identification of a separately defined zero-temperature posterior
is needed. If $\sigma$ denotes the strong-prior parameter in the convention
with density proportional to $\exp(-\kappa_N/(2\sigma^2))$, then
$t=\sqrt N\,\sigma$.

\begin{corollary}[Explicit width and depth powers]\label{cor:powers}
For fixed $d,c,\varepsilon,\delta$, put
\[
 F=1+\log^2(2+C)+\log(2+L),\qquad G=\log(2+CL).
\]
In either scale regime of Theorem~\ref{thm:main}, a sufficient sample
count has order
\begin{equation}\label{eq:joint-rate}
 O\!\left(C^4L^7F^4\log^2(2+L)G\right)
   =\widetilde O(C^4L^7).
\end{equation}
At fixed depth this is $O(C^4\log^9(2+C))$; at fixed width it is
$O(L^7\log^7(2+L))$.
\end{corollary}

The exact statistical dependence is in \eqref{eq:budget}--\eqref{eq:samples}.
In the second scale regime accuracy also enters $K$ through
$\log(2/\varepsilon)$. The asymptotic display \eqref{eq:joint-rate} holds
other resources fixed; it is not an optimality claim.

For the usual confidence criterion, define
$p_S(x)=Q_{S,N,t,c}\{f(x)z_w(x)>1/4\}$. Then
$D\{p_S\le2/3\}\le3R(Q)$; using $\varepsilon/3$ in the theorem gives
correct-margin posterior probability strictly greater than $2/3$ on all
but an $\varepsilon$ fraction of test inputs. This also gives correct-sign
confidence and covers arbitrary posterior mixtures.

\section{A quantized sampling construction}

We prove the deterministic finite-code statement needed for the posterior.
Fix $N$ and a network $w$ with $\kappa_N(w)\le K$. Write
$k_s=\norm{W^s}_F^2$ for $2\le s\le H$ and
$k_L=N\norm v_2^2$. Their sum is at most $K$.

\subsection{Pruning and conditional concentration}
In every internal row and in the readout, remove entries with absolute
value greater than
\[
 \frac{8\sqrt K\,\norm{\text{row}}_2}{r\sqrt N}.
\]
Zero rows remain zero. Denote pruned arrays by $\bar W^s$ and $\bar v$.
The absolute mass removed from a row is at most its squared norm divided
by this threshold. On any vector in $[-1,1]^N$, the internal pruning error
therefore has Euclidean norm at most
$r\sqrt{Nk_s/K}/8$; the readout error is at most
$r\sqrt{k_L/K}/8$.

Let $S^1,\ldots,S^H$ be independent lists of $m$ uniform indices in
$\{1,\ldots,N\}$, with replacement. Let $\mathcal Q_\zeta$ round a real
number to a nearest multiple of $\zeta$, with a fixed tie convention.
It preserves zero and has error at most $\zeta$. Define the \emph{full
virtual states} by
\begin{align}
 \widehat h^1_i(x)&=\tanh\!\left(
       \sum_{j=1}^{d+1}\mathcal Q_\zeta(W^1_{ij})\bar x_j\right),\notag\\
 \widehat h^s_i(x)&=\tanh\!\left(
       \sum_{j=1}^m
       \mathcal Q_\zeta\!\left(\frac Nm\bar W^s_{i,S^{s-1}_j}\right)
       \widehat h^{s-1}_{S^{s-1}_j}(x)\right),\quad 2\le s\le H,\notag\\
 \widehat z(x)&=\sum_{j=1}^m
       \mathcal Q_\zeta\!\left(\frac Nm\bar v_{S^H_j}\right)
       \widehat h^H_{S^H_j}(x).
 \label{eq:virtual}
\end{align}
The stored coordinates at the lists $S^s$ form an ordinary width-$m$
tanh network. This holds even when $m>N$ or indices repeat.

For each $s\ge2$, the state $\widehat h^{s-1}$ depends only on
$S^1,\ldots,S^{s-2}$. It remains in $[-1,1]^N$ and is independent of
the fresh list $S^{s-1}$. Conditional on the preceding lists, for one
uniform index $I$ define
\[
 Y=N\bar W^s_{:,I}\widehat h^{s-1}_I.
\]
The pruning threshold and bounded states give
\[
 \E\norm{Y-\E Y}_2^2\le Nk_s,\qquad
 \norm{Y-\E Y}_2\le16\sqrt{NKk_s}/r.
\]
Hilbert-space Bernstein \cite{tarres-yao}, applied at tolerance
$r\sqrt{Nk_s/K}/8$, bounds the failure probability by
\[
 2\exp\!\left(-\frac{3mr^2}{640K}\right)
 \le 2\exp\!\left(-\frac{mr^2}{256K}\right).
\]
The readout uses the independent final list $S^H$ and has the same
calculation without factors $\sqrt N$. Zero-energy stages are exact.
A union bound over the $2^d$ inputs and at most $L$ stages gives
simultaneous success probability at least $199/200$ by \eqref{eq:budget}.
The concentration estimate is applied \emph{before} coefficient rounding
at the current stage, but at the already rounded preceding virtual state.
This is why its conditional independence remains valid.

\subsection{All errors are disturbances of the original network}
On this concentration event, pruning and sampling together give stage
preactivation errors at most $r\sqrt{k_s/K}/4$ in normalized Euclidean
norm, with the scalar counterpart at the readout. Tanh is 1-Lipschitz,
so the corresponding stacked postactivation/readout error has norm at
most $r/4$.

Rounding adds at most $(d+1)\zeta$ per first-layer coordinate, and at
most $m\zeta$ per coordinate in every later layer and at the readout.
This uses only that all virtual hidden states in \eqref{eq:virtual} lie
in $[-1,1]^N$. The stacked rounding error thus satisfies
\[
 \norm{\eta_{\rm round}}_*^2
 \le\zeta^2\bigl[(d+1)^2+Hm^2\bigr]=r^2/16.
\]
Define the actual disturbances in the original, unpruned network by
\[
 \eta^1=\widehat h^1-\tanh(W^1\bar x),\qquad
 \eta^s=\widehat h^s-\tanh(W^s\widehat h^{s-1}),\qquad
 \eta^{\rm out}=\widehat z-v^\top\widehat h^H.
\]
They reproduce the virtual states and output exactly. Splitting each
into the pruning/sampling and rounding contributions and using the
triangle inequality in the stacked norm gives
\begin{equation}\label{eq:noise-budget}
 \norm\eta_*\le r/4+r/4=r/2.
\end{equation}
If $w\in\mathcal R_{N,c}$, then $r\le r_c(w)$ because $E_v(w)\le K$.
The original robustness assumption therefore yields
\begin{equation}\label{eq:approximation}
 \sup_{x\in\mathcal X}|T(z_w(x))-T(\widehat z(x))|\le1/16.
\end{equation}
There is no subsequent unprotected rounding step, and no assertion that
the sampled network itself satisfies \eqref{eq:support}.

\subsection{A finite code independent of learner width}
Consider the same sampled coefficients before applying
$\mathcal Q_\zeta$. Their energy expectations obey
\[
 \frac dm\E\norm{W^1_{\rm comp}}_F^2
   =\frac dN\norm{W^1}_F^2,\qquad
 \E\norm{W^s_{\rm comp}}_F^2\le\norm{W^s}_F^2,\qquad
 m\E\norm{v_{\rm comp}}_2^2\le N\norm v_2^2.
\]
These statements concern sampled parameters, not virtual activations.
They use independence of adjacent index lists and are unaffected by
rounding inside \eqref{eq:virtual}. Markov's inequality gives sampled
unrounded energy at most $4K$ with probability at least $3/4$.
This event intersects the concentration event.

On the intersection, every unrounded stored coefficient has absolute value
at most $B=\sqrt{4mK}$. Its rounded value belongs to a fixed grid containing
at most $3+2B/\zeta$ points. Although rounding was defined on all real
numbers for the full virtual computation, only these bounded stored
coefficients enter its width-$m$ realization. Thus a fixed class
$\mathcal G$ of rounded networks satisfies
\begin{equation}\label{eq:code}
 \log|\mathcal G|\le P\log(3+2B/\zeta)=J,
\end{equation}
and every internally robust energy-$K$ network at every width has a member
of this class satisfying \eqref{eq:approximation}.

\section{Posterior learning and the sample polynomial}

\begin{proof}[Proof of Theorem~\ref{thm:main}]
Let $\psi$ equal one on $(-\infty,1/3]$, zero on $[5/12,\infty)$,
and be affine between these thresholds. Every robust network of energy
at most $K$ has a rounded representative $g\in\mathcal G$.
If it fits a training margin, then
\[
 f(X_i)T(g(X_i))\ge1/2-1/16=7/16>5/12.
\]
Every test input where $f(x)z_w(x)\le1/4$ satisfies
$f(x)T(g(x))\le1/4+1/16=5/16<1/3$.
The representative therefore has zero empirical ramp loss, and its
population ramp loss bounds the original network's test-margin error.

For a fixed nonnegative loss bounded by one with population mean $u$,
the chance that all $n$ iid losses equal zero is at most
$(1-u)^n\le e^{-nu}$. A union bound using \eqref{eq:code} and
\eqref{eq:samples} shows that, with probability at least $1-\delta$,
every zero-empirical-loss representative has population loss at most
$\varepsilon/2$. This single sample event controls every robust fitting
network of energy at most $K$, uniformly over width and scale.

The target cylinder bounds the raw conditioning probability below by
$e^{-ND_*/t^2}$ for every correctly labeled sample.
Under the raw Gaussian, $N\kappa_N/t^2$ is chi-square with $P_N$
degrees of freedom. Its exponential moment gives
\begin{equation}\label{eq:tail}
 Q_{S,N,t,c}\{\kappa_N>K\}
 \le\exp\!\left(-\frac{NK}{4t^2}
                  +\frac{P_N}{2}\log2+\frac{ND_*}{t^2}\right).
\end{equation}
For $K=8D_*$, the first scale threshold makes this at most
$\varepsilon/2$. For the second choice of $K$, use
$t^2P_N/N\le d+L$ when $t^2\le1/N$. The exponent multiplied by
$t^2/N$ is at most
\[
 -K/4+D_*+(d+L)(\log2)/2\le-\log(2/\varepsilon).
\]
Since $N/t^2\ge1$, the tail is again at most $\varepsilon/2$.
Adding this discarded posterior mass to the retained risk proves the
claim for arbitrary posterior mixtures.
\end{proof}

\begin{proof}[Proof of Corollary~\ref{cor:powers}]
Put $T_L=d+\log(400L)$. The budget gives
$m=O(Kr^{-2}T_L)$. All factors in $B/\zeta$ are polynomial in
$d,L,K,r^{-1}$, so
\[
 J=O\!\left(
 [dKr^{-2}T_L+LK^2r^{-4}T_L^2]
 \log(2+d+L+K+r^{-1})\right).
\]
The target constants imply
\[
 D_*=O\!\left((d+LC)
 [\log^2(2C)+\log(2L)+\log(2+d+C)]+(1+c^2)C\right).
\]
At fixed $d,c,\varepsilon,\delta$, it follows that
\[
 K=O(CLF),\quad r^{-2}=O(CL^2F),\quad
 m=O(C^2L^3F^2\log(2+L)),
\]
\[
 P=O(C^4L^7F^4\log^2(2+L)),\qquad
 \log(3+2B/\zeta)=O(G).
\]
Substitution proves \eqref{eq:joint-rate} and its fixed-width and
fixed-depth consequences.
\end{proof}

\section{Scope of the two compression arguments}

The energy-based baseline just proved has a total-energy cutoff.
At mean-field scaling, typical internal Frobenius energy grows with
learner width, so that particular proof cannot give its stated quartic
rate throughout the scale range. Entropy-metric compression with Gaussian-weighted coding avoids a
width-independent total-energy cutoff and proves the all-scale theorem
instead. Both methods use the same prior support.

The new depth-six tradeoff and all-scale quartic bound are sufficient
rates. Quadratic teacher-width dependence under this unchanged support,
and a depth-linear bound with quartic width, remain open. No matching
lower bound for the displayed polynomial exponents is asserted.
The proofs build on the target-cylinder and reconstruction estimates
of~\cite{robust-paper,quartic,all-scales}; their nontrivial imported
lemmas are stated explicitly below. These new arguments and rates are
not Lean formalized. Existing formal proofs for earlier bounds do not
certify them.

\appendix
\section{A robust target cylinder at every width}
\label{app:cylinder}

We give the probability bound used in the statistical proof. It uses a
robust digital implementation, not a claim that every implementation is
stable. Let $N\ge C$. Partition the neurons into $C$ balanced clone groups
of sizes $k_a\in\{\lfloor N/C\rfloor,\lceil N/C\rceil\}$, so
$C/2\le N/k_a\le2C$. Let $P$ project onto vectors constant on each clone
group. Fill unused logical slots with available constants or copies of an
existing Boolean wire. The ideal wire
vector $b^s(x)$ then has entries in $\{-1,1\}$, belongs to the range of
$P$, and has Euclidean norm $\sqrt N$.

Use gain $A$ in the tanh simulation of every Boolean gate. For example,
NAND on $u,v\in\{-1,1\}$ is the sign of $1-u-v$; copies and negations
use $u$ and $-u$. Bias constants in internal layers are supplied by counted
constant wires. The ideal signed preactivation of every gate is at least
$A$. In the clone implementation, each coefficient from a logical wire is
divided equally among that wire's clones. Set
\begin{align*}
 e_*&=\frac1{16\sqrt{2C}},&
 \tau_0&=\frac{e_*}{4},&
 b_L&=\sqrt{16\log9+8\log(4L)},\\
 A&=\max\{32b_L,2\log(8/e_*)\},&
 a&=\max\{2,256c\},&
 \rho&=\min\{1/4,A/[4(d+1)],Ae_*/(64\sqrt C)\}.
\end{align*}
The teacher readout averages the clones of the output wire with total
coefficient $a$.

In Gaussian coordinates $\sqrt d\,W^1$, $\sqrt N\,W^s$, and $Nv$, all
entries have variance $t^2$. Restrict each first-layer coordinate, each
internal row's projection onto every normalized clone-group direction, and
each readout coordinate to an interval of radius $\rho t$ about its teacher
value. The internal matrices then decompose as
\[
 W^s=W_*^s+R^s+E^s,\qquad
 R^s=R^sP,\quad \|R^s\|_{\rm op}\le\rho t\sqrt C,
 \quad E^sP=0.
\]
The remaining Gaussian bulk $E^s$ is independent of the constrained
coordinates. A $1/4$-net of the unit sphere has at most $9^N$ points,
and the usual two-net Gaussian bound gives
\[
 \Pr\{\|E^s\|_{\rm op}>b_Lt\}
 \le2\exp(2N\log9-Nb_L^2/8)\le\frac1{2L}.
\]
Thus all bulk norms are at most $b_Lt$ with conditional probability at
least $1/2$.

\paragraph{Finite disturbance robustness.}
Allow a separate RMS disturbance of size $\tau_0$ at each hidden layer and
scalar output disturbance at most $\tau_0$. This includes the entire
total-budget ball $\|\eta\|_*\le\tau_0$.
Write $e_s=\|h^s_\eta-b^s\|_2/\sqrt N$.
If $e_{s-1}\le e_*$, every clone-group mean differs from its ideal bit by
at most $e_*\sqrt{N/k_a}\le1/16$. At most three teacher coefficients,
including a constant coefficient, enter a gate, so its intended signed
preactivation remains at least $A/2$.
For $\epsilon\in\{-1,1\}$ and real $u$, the elementary global secant bound
\[
 \epsilon q\ge A/2
 \quad\Longrightarrow\quad
 |\tanh(q+u)-\epsilon|\le2e^{-A/2}+8|u|/A
\]
holds: for $|u|<A/4$ use saturation, and otherwise the linear term is at
least two. Since $E^sb^{s-1}=0$ and
$\|h^{s-1}_\eta\|_2/\sqrt N\le1+\tau_0$, it follows that
\[
 e_s\le2e^{-A/2}+
          (8/A)\rho\sqrt C(1+\tau_0)
          +(8b_L/A)e_{s-1}+\tau_0\le e_*.
\]
Each of the four terms is at most $e_*/4$.
The first layer initializes the same invariant: its coordinatewise
preactivation perturbation is at most $\rho(d+1)\le A/4$, and its
postactivation disturbance contributes at most $\tau_0$ in RMS norm.
Consequently the noisy signed readout is at least
\[
 f(x)z_w(x;\eta)
 \ge\frac{15a}{16}-\rho(1+\tau_0)-\tau_0>1
 \qquad(x\in\mathcal X,\ \|\eta\|_*\le\tau_0).
\]

The readout center satisfies $\sqrt N\|v_*\|\ge a\sqrt{C/2}$, and its
permitted perturbation has $\sqrt N\|v-v_*\|\le\rho$.
Therefore $E_v\ge a^2C/8$, whence
\[
 r_c(w)\le\frac{2\sqrt2c}{a\sqrt C}
       \le\frac1{64\sqrt{2C}}=\tau_0.
\]
All clean and noisy clipped scores on the calibrated ball equal $f(x)/2$.
The entire cylinder, with the bulk-norm event imposed, therefore belongs
to $\mathcal R_{N,c}$ and satisfies the training likelihood for every sample.
This also proves why every fixed $c>0$ is allowed: a larger calibration
requires only a polynomially larger readout coefficient.

\paragraph{Gaussian mass.}
There are $Nq$ constrained independent scalar coordinates, where
$q=d+2+(L-2)C$. The sum of their squared centers is at most $NB_0$, with
\[
 B_0=9dA^2+18(L-2)CA^2+2a^2C.
\]
These loose constants cover the fan-in-two gate implementation above.
For $Z\sim\mathcal N(0,t^2)$, $\rho\le1/4$, and any center $m$,
\[
 \Pr\{|Z-m|\le\rho t\}
 \ge\frac{2\rho}{\exp(1)\sqrt{2\pi}}\exp(-m^2/t^2).
\]
Multiplication over the coordinates and the independent bulk event gives
\begin{equation}\label{eq:standalone-cylinder-mass}
 \mu_{N,t}(\text{target cylinder})
 \ge\exp(-ND_*/t^2),\qquad
 D_*=B_0+q\log\frac{\exp(1)\sqrt{2\pi}}{2\rho}+\log2.
\end{equation}
Here $D_*\ge1$ and is bounded by a universal polynomial in $d,L,C,c$.
The construction and estimate hold at every integer width $N\ge C$.

\section{Sequential quantization at every mean-field-or-stronger scale}
\label{app:all-scales-quantization}

The energy-based baseline compressor requires a bound
on total network energy. At the finite mean-field scale, the unconditioned
Gaussian background can have growing total energy. An entropy-based
compressor avoids that difficulty. In this appendix we combine the same
sequential-quantization principle with the active-direction and covariance
reconstruction estimates of~\cite{all-scales}. This supporting code gives a sufficient bound at every $0<t\le1$,
including $\gamma=1$; Appendix~\ref{app:weighted-code} gives the
sharper headline bound.
This appendix is not a new proof of those active-direction estimates: we
state precisely which properties are imported and prove the modification
that removes the coding precision's extra depth factor. Neither this
modification nor its new depth exponent is Lean formalized.

\subsection{Budgets and learning statement}

Use the model, support $\mathcal R_{N,c}$, and posterior
$Q_{S,N,t,c}$ from \eqref{eq:support}--\eqref{eq:conditioning}.
Let $D_*$ be any target-cylinder cost satisfying the conclusion of
Appendix~\ref{app:cylinder}. Thus the raw posterior conditioning event
has probability at least $\exp(-ND_*/t^2)$, for every correctly labeled
sample and every $N\ge C$, $0<t\le1$. For
$0<\varepsilon,\delta<1$, put
\begin{align*}
 \beta&=\varepsilon/4,&
 B_v&=\max\{8,8[D_*+\log(1/\beta)]\},&
 K&=D_*+1,\\
 r&=\min\{1/6,c/\sqrt{1+(L-1)B_v}\},&
 \alpha&=\varepsilon/12.
\end{align*}
Define the compression parameters
\begin{align*}
 e&=\frac{\alpha r^2}{1000L},&
 \eta&=\frac{\alpha}{100L},&
 q_a&=\left\lceil\frac{2K}{e\eta}\right\rceil,\\
 D_0&=[(d+1)q_a+2]\log2,&
 V&=8K+2q_a,&
 B_H^2&=4V/e,\\
 K_0&=(4K+2\log2)/\eta,&
 B_1^2&=4[4K+2(d+1)\log2]/e.
\end{align*}
The letter $q_a$ denotes the number of retained active directions; it
is unrelated to a population partition or a support restriction.
Let $m$ be the least integer no smaller than each of
\begin{align*}
 &10000(V+q_a+1)/e,\\
 &64(q_a+1)\log(4(q_a+1)/\eta),\\
 &512K_0r^{-2}[d\log2+\log(2/\eta)+1].
\end{align*}
Finally set
\begin{align*}
 M&=\max\{m,\lceil LD_0/K\rceil,1\},&
 B_G&=\max\{2,2\sqrt{\log(6m/e)}\},\\
 P_*&=M(d+1)+(L-2)M^2+M,&
 B_{\rm small}^2&=\frac{4(4MK+2P_*\log2)}\alpha,\\
 B_{\rm large}&=\max\left\{
 B_1,\frac{2B_H\sqrt{q_a+1}}m+\frac{\sqrt2 B_G}{\sqrt m},
 \frac{\sqrt{K_0}}m,1\right\},\\
 B_*&=1+\max\{B_{\rm small},B_{\rm large},1\},&
 \zeta&=\frac{\sqrt e}{4(M+d+1)},\\
 J_*&=P_*\log(3+4B_*/\zeta).
\end{align*}
All square roots above denote the nonnegative root.

\begin{theorem}[An improved all-scale budget]
\label{thm:all-scales-sequential}
Under the Boolean-teacher convention of the main theorem, the condition
\begin{equation}\label{eq:all-scales-sequential-budget}
 n\ge\frac{16}{\varepsilon}\left[J_*+\log(1/\delta)\right]
\end{equation}
implies, with probability at least $1-\delta$ over iid training inputs,
\[
 \E_{w\sim Q_{S,N,t,c},\,X\sim D}
     \mathbf1\{f(X)z_w(X)\le1/4\}\le\varepsilon
\]
simultaneously for every $N\ge C$ and every $0<t\le1$.
The conclusion is also valid if global clean-output-margin conditioning
is added to the prior. No population-merging or additional robustness
assumption is made.
\end{theorem}

With $d,c,\varepsilon,\delta$ fixed, write
\[
 F=1+\log^2(2+C)+\log(2+L),\qquad G=\log(2+CL).
\]
The target constants and the budgets above give
\begin{equation}\label{eq:all-scales-sequential-rate}
 n=O(C^6L^{17}F^6G).
\end{equation}
This improves the former $O(C^6L^{18}F^6G)$ budget while preserving its
fixed-depth $O(C^6\log^{13}(2+C))$ width consequence.
A further target-cylinder sharpening below removes the $\log(2+L)$
term from $F$.

\subsection{Imported entropy and reconstruction estimates}

We use the following properties of the construction in the appendix
``Finite posterior compression by internal-noise coupling''
of~\cite{all-scales}. Its lemmas on future-weight entropy, whole-layer
active directions, and whitened sampling prove these properties for
\emph{any bounded virtual prefix}. They do not require a Gaussian law
for the prefix and do not assume its robustness.

First, suppose a prefix is $A$, the current original weight block is $X$,
and the future blocks are $Y$, with product Gaussian reference law
$P_X\otimes P_Y$. For a finite message $J$, a retained quantity
$H_0=f(A,J,X)$, and a new prefix $A'$ which is a randomized function of
$(A,J,H_0)$ using fresh randomness conditionally independent of $Y$,
\begin{equation}\label{eq:imported-future-entropy}
 \E_{A'}D_{\rm KL}(Q_{Y\mid A'}\Vert P_Y)
 \le \E_A D_{\rm KL}(Q_{X,Y\mid A}\Vert P_X\otimes P_Y)
       +\mathcal H(J\mid A).
\end{equation}
Here $\mathcal H$ denotes conditional Shannon entropy. The inequality
is the KL chain rule and data processing; deterministic processing of
$A'$ cannot increase its left side.

Second, if $Q\ll\mu_{N,t}$ has
$D_{\rm KL}(Q\Vert\mu_{N,t})\le NK/t^2$ and $N>M$, the budgets above
permit the following sequential construction. Its current virtual layer
has $N$ rows, each evaluated on $m$ stored source neurons from the
preceding layer. The first original layer is clipped rowwise at
physical standardized norm $B_1$. Each internal stage retains at most
$q_a$ active input directions, clips the active row norms at $B_H$,
reconstructs the complementary covariance by whitened source sampling,
and clips its new Gaussian coefficients at $B_G$. The finite message
selecting active inputs has entropy at most $D_0$; all other new
randomness is fresh conditional on the prefix and active matrix. Hence
\eqref{eq:imported-future-entropy} bounds the unprocessed original-weight
entropy at stage $s$ by
\[
 B_s\le NK/t^2+sD_0\le2NK/t^2.
\]
For every fixed input $x$, the conditional reconstruction couplings give
expected squared RMS activation injection less than $20e$ per hidden
stage, with success indicators. The first layer has the smaller bound
$e$. Before any later row selection, each new internal coefficient has
absolute value at most
\[
 \frac{2B_H\sqrt{q_a+1}}m+\frac{\sqrt2B_G}{\sqrt m}.
\]
The readout is truncated at $N\norm v^2\le K_0$ and reconstructed using
$m$ source samples, with coefficient bound $\sqrt{K_0}/m$ and uniform
output error at most $r/16$. The sum of probabilities of rank, Gram,
and readout failures is at most $3L\eta$.

For clarity, the quantitative source-reconstruction input is the
following estimate. For bounded source activations $|h_j(x)|\le1$,
active dimension at most $q_a$, and active matrix $H_0$, its whitened
half-uniform, half-leverage sampling law has expected squared RMS
error, with the good-Gram indicator, at most
\[
 \frac{2\norm{H_0}_F^2/N+16(q_a+1)}m.
\]
The averaged active energy is bounded by $V$, while the Gram failure
probability is at most $\eta$. These facts, Gaussian clipping, the
active-direction transport estimate, and the squared triangle
inequality give the stated $20e$ bound. Each bound holds for every
bounded prefix to which the entropy estimate applies. These are the
nontrivial reconstruction estimates imported from~\cite{all-scales};
the argument below does not infer them merely from an existing
compressed network's clean-output accuracy.

\subsection{Rounding inside the canonical construction}

Choose fixed scalar lattice rounding with error at most $\zeta$ and
with zero fixed. After first-layer clipping, round all its coefficients
before constructing any later prefix. At each internal stage apply the
imported reconstruction procedure to the \emph{already rounded} prefix,
then round every newly constructed virtual row before processing the
next original layer. Finally round the reconstructed readout. Keep the
source rows selected at subsequent stages, as in the unrounded
construction. There is no terminal rounding pass.

This modification preserves the preceding estimates. Tanh keeps the
rounded virtual activations in $[-1,1]$. Each preceding state is
specified before the next fresh source list is sampled. Thus the
boundedness and conditional-independence assumptions of the sampling
estimate remain valid. Rounding new rows is a deterministic function
of the same prefix, message, active matrix and fresh Gaussian variables
as their unrounded construction. It therefore preserves
\eqref{eq:imported-future-entropy} without adding a message.
The resulting prefix may have a discrete law; that law never receives
a Gaussian relative-entropy comparison.

Fix an input $x$. The conditional transports can be glued while keeping
the original full-network marginal equal to $Q$ and the outgoing
rounded-prefix/future-weight marginal equal to the next canonical
stage's law. To do so, extend each transition by the previous
conditional law of past original weights given its incoming prefix and
remaining original weights. This disintegration preserves all earlier
computed prefixes. Define the new postactivation injections by
\[
 \eta^s=h^s_{\rm rounded\ virtual}(x)
       -\tanh\bigl(W^s_{\rm original}
                         h^{s-1}_{\rm rounded\ virtual}(x)\bigr).
\]
The final candidate output is exactly the original network's output
under these injections and the analogous readout injection. Its
marginal law is independent of $x$, although the coupling may depend
on $x$.

Coefficient rounding changes a first-layer preactivation by at most
$(d+1)\zeta$, and every later virtual preactivation by at most
$m\zeta$. Both are at most $\sqrt e/4$. Consequently the expected
squared RMS hidden injection, including rounding, is at most
\[
 2(20e)+2(e/16)<41e.
\]
Outside the abort events, Markov's inequality gives
\[
 \Prb\left\{\sum_{s<L}\norm{\eta^s}^2/N>r^2/16\right\}
 \le\frac{656Le}{r^2}=0.656\alpha.
\]
The total abort probability is at most $3L\eta=0.03\alpha$.
The readout rounding error is at most $m\zeta\le\sqrt e/4\le r/64$,
where the last inequality uses $\alpha\le1/8$ and $L\ge2$.
Together with its sampling error, the scalar injection is at most
$5r/64$. On the good event,
\[
 \norm\eta_*\le\sqrt{r^2/16+25r^2/4096}
       =\frac{\sqrt{281}}{64}r<r/2.
\]
The exceptional probability is less than $\alpha$.

When $N\le M$, the Gaussian entropy-moment inequality instead gives,
with probability at least $1-\alpha/4$, the coefficient bound
$B_{\rm small}$ above. On this event round every original coefficient
and compute the rounded network sequentially; on its complement use a
zero-readout default. At each input the rounded computation is the
original network under injections satisfying
\[
 \norm\eta_*^2
 \le\zeta^2[(d+1)^2+(L-1)N^2]
 \le\zeta^2 L(M+d+1)^2=Le/16<r^2/4.
\]
This identity uses bounded rounded hidden activations, not a bound on
products of original matrix norms.

All retained coefficients, after rounding, lie in $[-B_*,B_*]$.
Padding each candidate to width $M$ with zero neurons is exact for tanh.
A grid with at most $3+4B_*/\zeta$ entries per coefficient gives a finite
class $\mathcal H_*$ with
\begin{equation}\label{eq:all-scales-code}
 \log|\mathcal H_*|\le J_*.
\end{equation}
The construction proves the following strengthened finite-code statement:
for every such $Q$ there is one law $\Gamma_Q$ on $\mathcal H_*$ such
that, at each fixed input, a coupling has
\[
 h(x)=z_w(x;\eta),\qquad \norm\eta_*<r/2
\]
except on probability at most $\alpha$. Only the coupling, not its two
marginals, can depend on the input.

\subsection{Posterior and statistical argument}

Under the raw Gaussian, $NE_v/t^2$ is chi-square with $N$ degrees of
freedom. The target-cylinder denominator and the exponential moment
bound give posterior mass at most $\beta$ to $E_v>B_v$. The retained
conditional posterior $Q'$ satisfies
\[
 D_{\rm KL}(Q'\Vert\mu_{N,t})
 \le ND_*/t^2+\log(1/(1-\beta))\le NK/t^2.
\]
Every retained original network is robust on the common radius $r$.
Apply the strengthened finite-code statement to $Q'$. On each good
coupling the clipped-score error is at most $1/16$.

Let $\psi$ be one on $(-\infty,1/3]$, zero on $[5/12,\infty)$,
and affine between these intervals. A correctly fitted training point
has candidate ramp loss zero on the good coupling. A test-margin
failure of the original network has candidate ramp loss one there.
For the finite class $\mathcal H_*$, a relative Chernoff inequality
and a union bound give, on a sample event of probability at least
$1-\delta$,
\[
 \E_D\psi(f(X)T(h(X)))
 \le \frac2n\sum_{i=1}^n\psi(f(X_i)T(h(X_i)))
       +\frac{8[J_*+\log(1/\delta)]}{n}
 \qquad(h\in\mathcal H_*).
\]
Averaging over $\Gamma_{Q'}$, comparing population risks, and restoring
the discarded posterior mass bounds the desired risk by
\[
 \beta+3\alpha+\frac{8[J_*+\log(1/\delta)]}{n}
 \le\varepsilon
\]
under \eqref{eq:all-scales-sequential-budget}. The finite class and sample
event do not depend on $N,t$. This proves
Theorem~\ref{thm:all-scales-sequential}. Adding global-margin conditioning
changes none of these steps, since the same target cylinder is retained.

\subsection{Width--depth accounting and a sharper target cylinder}

For fixed $d,c,\varepsilon,\delta$, the displayed budgets give
\begin{align*}
 D_*,K,B_v&=O(CLF),&r^{-2}&=O(CL^2F),\\
 e^{-1}&=O(CL^3F),&\eta^{-1}&=O(L),\\
 q_a,V&=O(C^2L^5F^2),&M&=O(C^3L^8F^3),\\
 P_*&=O(C^6L^{17}F^6).
\end{align*}
Both $B_*$ and $\zeta^{-1}$ are polynomial in $C,L,F$, so
$\log(3+4B_*/\zeta)=O(G)$. The old terminal grid had an additional
factor $L$ here. This establishes \eqref{eq:all-scales-sequential-rate}.

There is a separate improvement to the target-cylinder budget which
changes neither the prior nor the compression. In
Appendix~\ref{app:cylinder}, replace the depth-dependent bulk cutoff by
\[
 b=\sqrt{16\log9+8\log4},\qquad
 A=\max\{32b,2\log(8/e_*)\}.
\]
Each unconstrained Gaussian bulk matrix has, by the same two-net bound,
\[
 \Prb\{\norm{E^s}_{\rm op}>bt\}
 \le2\exp(2N\log9-Nb^2/8)=2\cdot4^{-N}\le1/2.
\]
The bulk matrices are conditionally independent across layers.
Their simultaneous success probability is therefore at least
$2^{-(H-1)}$. Every deterministic stability estimate in the cylinder
proof is unchanged. Recompute $B_0,q_0,\rho$ with this constant-cutoff
choice of $A$ and use the conservative cost
\begin{equation}\label{eq:product-cylinder-cost}
 D_*^{\rm prod}=B_0+q_0\log\frac{\exp(1)\sqrt{2\pi}}{2\rho}+H\log2.
\end{equation}
Indeed $N/t^2\ge1$ absorbs the product-event cost into this rate.
At fixed $d,c$, with $F_0=1+\log^2(2+C)$, one has
$D_*^{\rm prod}=O(CLF_0)$. Substituting this valid cylinder cost in
Theorem~\ref{thm:all-scales-sequential} gives
\begin{corollary}[Depth logarithms removed from the target gain]
\label{cor:all-scales-product-cylinder}
For fixed $d,c,\varepsilon,\delta$, the unchanged robust prior admits
sufficient sample count
\[
 O\!\left(C^6L^{17}[1+\log^2(2+C)]^6\log(2+CL)\right)
\]
at every $N\ge C$ and every $0<t\le1$. At fixed teacher width this is
$O(L^{17}\log(2+L))$.
\end{corollary}
The same cylinder replacement is also available in the quartic main
theorem; it is independent of the sampling and quantization argument.
Its proof uses the target constants only through their valid mass lower
bound. With this substitution the quartic sufficient order becomes
\[
 O\!\left(C^4L^7[1+\log^2(2+C)]^4\log^2(2+L)\log(2+CL)\right).
\]
At fixed teacher width this is $O(L^7\log^3(2+L))$.

\begin{remark}[A conservative joint polynomial]
Put $S=2+d+L+C+c+c^{-1}+\varepsilon^{-1}+\log(1/\delta)$.
The coarse accounting in~\cite{all-scales} for the unchanged budgets gives
$e^{-1}=O(S^8)$, $M=O(S^{21})$, $P_*=O(S^{43})$, and
$B_*=O(S^{22})$. The new mesh has
$\zeta^{-1}=O(S^{25})$. Thus $J_*=O(S^{43}\log S)$ and
\eqref{eq:all-scales-sequential-budget} is $O(S^{44}\log S)$,
hence $O(S^{45})$. This deliberately loose degree-45 envelope is a
consequence of the conventional modification proved here. It does not
replace or enlarge the existing Lean certification of the earlier
degree-46 envelope.
\end{remark}

\section{Bounded-density compression with factored Gaussian reconstruction}
\label{app:factor-code}

The all-scale code in Appendix~\ref{app:all-scales-quantization} can be
sharpened in two further ways. First, the actual conditioned posterior has
a bounded density relative to the raw Gaussian. Propagating this stronger
property through the virtual computation removes the factor $1/\eta$
from the active-direction cutoff. This argument is specific to bounded
density; it does not upgrade the compression theorem for arbitrary laws
having only a relative-entropy bound. Second, the active part of each
reconstructed matrix already has small rank, and its fresh Gaussian part
can be replaced by a small-rank Gaussian factorization. Coding the factors
then gives a sufficient sample order $\widetilde O(C^5L^{12})$ at every
$0<t\le1$. The support remains exactly $\mathcal R_{N,c}$. This is a supporting
code: the entropy-metric and Gaussian-weighted constructions in the
following appendices improve its fifth-power width dependence to the
headline quartic bound.

We give the density and factorization arguments below. We use the
whole-layer active-direction and whitened-sampling estimates stated in
Appendix~\ref{app:all-scales-quantization}, with their original proofs
in~\cite{all-scales}. As in that appendix, every comparison coupling is
allowed to depend on the input, whereas the candidate law is not. The
argument and the improved exponents are conventional proofs and are not
Lean formalized.

\subsection{Explicit budgets and learning statement}

Use the product-cylinder cost $D_*=D_*^{\rm prod}$ from
\eqref{eq:product-cylinder-cost}. For $0<\varepsilon,\delta<1$, define
\begin{align}
 \beta&=\varepsilon/4,&
 B_v&=\max\{8,8[D_*+\log(1/\beta)]\},&
 K&=D_*+1,\notag\\
 r&=\min\{1/6,c/\sqrt{1+(L-1)B_v}\},&
 \alpha&=\varepsilon/12,&
 e&=\frac{\alpha r^2}{1000L},\qquad
 \eta=\frac{\alpha}{100L}.
 \label{eq:factor-posterior-budgets}
\end{align}
Set
\begin{align}
 q&=\left\lceil\frac{2K+\log(1/\eta)}e\right\rceil,&
 D_0&=[(d+1)q+2]\log2,&
 V&=8K+2q,\notag\\
 B_H^2&=4V/e,&
 K_0&=(4K+2\log2)/\eta,&
 B_1^2&=4[4K+2(d+1)\log2]/e.
 \label{eq:factor-active-budgets}
\end{align}
Let $m$ be the least integer no smaller than each of
\begin{align}
 &10000(V+q+1)/e,\notag\\
 &64(q+1)\log(4(q+1)/\eta),\label{eq:factor-source-width}\\
 &512K_0r^{-2}[d\log2+\log(2/\eta)+1].\notag
\end{align}
The factor rank and the small-width cutoff are
\begin{equation}\label{eq:factor-rank-cutoff}
 p=\lceil1000/e\rceil,\qquad k=q+p,\qquad
 n_0=\max\{1,\lceil LD_0/K\rceil\}.
\end{equation}
In particular, the cutoff $n_0$ need not be as large as $m$.
Define factor bounds and parameter counts by
\begin{align}
 B_g^2&=4p/e,\qquad A_f=\max\{B_H,B_g,1\},\notag\\
 B_f&=\max\{2\sqrt{q+1}/m,\sqrt{2/\eta},1\},\notag\\
 P_{\rm small}&=n_0(d+1)+(L-2)n_0^2+n_0,\notag\\
 P_{\rm fac}&=m(d+1)+2(L-2)mk+m,\notag\\
 B_{\rm small}^2&=\frac{4(4n_0K+2P_{\rm small}\log2)}\alpha,\notag\\
 B_{\rm code}&=1+\max\{B_{\rm small},B_1,A_f,B_f,
                                      \sqrt{K_0}/m,1\}.
 \label{eq:factor-code-bounds}
\end{align}
Finally, choose the common scalar mesh
\begin{equation}\label{eq:factor-mesh}
 \zeta=\min\left\{1,
 \frac{\sqrt e}{4(d+1)},\frac{\sqrt e}{4m},
 \frac{\sqrt e}{4mk(A_f+B_f+1)},
 \frac{\sqrt e}{4(n_0+d+1)}\right\},
\end{equation}
and put
\begin{equation}\label{eq:factor-code-budget}
 J_{\rm fac}=\log2+
 \max\{P_{\rm small},P_{\rm fac}\}
               \log(3+4B_{\rm code}/\zeta).
\end{equation}
All these quantities are independent of learner width and Gaussian scale.

\begin{theorem}[All-scale learning by a factored code]
\label{thm:factor-all-scales}
Under the Boolean-teacher convention of the main theorem, the condition
\begin{equation}\label{eq:factor-sample-budget}
 n\ge\frac{16}{\varepsilon}
             [J_{\rm fac}+\log(1/\delta)]
\end{equation}
implies, with probability at least $1-\delta$ over iid training inputs,
\[
 \E_{w\sim Q_{S,N,t,c},\,X\sim D}
       \mathbf1\{f(X)z_w(X)\le1/4\}\le\varepsilon
\]
simultaneously for every $N\ge C$ and every $0<t\le1$.
The same conclusion holds if the posterior conditioning additionally
requires the label-independent global margin $|z_w(x)|>1/2$ for every
input $x$.
No population-merging assumption or additional internal-noise condition
is imposed.
\end{theorem}

\begin{corollary}[Improved width and depth powers]
\label{cor:factor-powers}
For fixed $d,c,\varepsilon,\delta$, let $F_0=1+\log^2(2+C)$.
A sufficient sample count in Theorem~\ref{thm:factor-all-scales} has order
\begin{equation}\label{eq:factor-powers}
 O\!\left(C^5L^{12}F_0^5\log(2+CL)\right)
       =\widetilde O(C^5L^{12}).
\end{equation}
At fixed teacher width this is $O(L^{12}\log(2+L))$; at fixed depth it is
$O(C^5\log^{11}(2+C))$.
\end{corollary}

\subsection{A bounded-density invariant for unprocessed weights}

The following lemma strengthens the entropy bookkeeping at the point
where an active rank must hold with high probability. A measure inequality
means domination on every measurable set.

\begin{lemma}[Conditional density and finite-message propagation]
\label{lem:factor-density}
Suppose that
\[
 Q_{A,X,Y}\le e^b(\nu_A\otimes P_X\otimes P_Y),
\]
where $\nu_A,P_X,P_Y$ are probability reference laws. Write
$u=dQ_{A,Y}/d(\nu_A\otimes P_Y)$. Then
\[
 Q\{u(A,Y)<\eta\}\le\eta,
\]
and, outside this event,
\begin{equation}\label{eq:factor-conditional-entropy}
 D_{\rm KL}(Q_{X\mid A,Y}\Vert P_X)
       \le b+\log(1/\eta).
\end{equation}
Suppose also that $J=j(A,Y)$ takes values in a fixed finite alphabet
$\mathcal J$, that $H_0=h(A,J,X)$, and that a new prefix $A'$ is generated
from $(A,J,H_0)$ through a kernel whose fresh randomness has no further
dependence on $Y$. There is a probability reference law $\nu_{A'}$ such
that
\begin{equation}\label{eq:factor-density-propagation}
 Q_{A',Y}\le e^{b+\log|\mathcal J|}
                            (\nu_{A'}\otimes P_Y).
\end{equation}
The original domination also implies
\[
 \E_{Q_A}D_{\rm KL}(Q_{X,Y\mid A}\Vert P_X\otimes P_Y)\le b.
\]
\end{lemma}
\begin{proof}
Let $F\le e^b$ be the joint density relative to
$\nu_A\otimes P_X\otimes P_Y$. Disintegration gives conditional density
$F(a,x,y)/u(a,y)$ for $Q_{X\mid a,y}$ wherever $u>0$.
Thus its relative entropy is at most $b-\log u(a,y)$.
Also
\[
 Q\{u<\eta\}=
 \int_{\{u<\eta\}}u\,d(\nu_A\otimes P_Y)\le\eta.
\]
This is an unconditional probability bound under the canonical joint law;
it need not hold separately at every fixed prefix.

For propagation, augment the reference law by an independent uniformly
distributed message in $\mathcal J$. The actual augmented law satisfies
\[
 Q_{A,J,X,Y}\le e^b|\mathcal J|
  (\nu_A\otimes\operatorname{Unif}(\mathcal J)
                       \otimes P_X\otimes P_Y).
\]
Apply the retained-data map and the same prefix kernel to both sides.
Under the reference law, its inputs and fresh randomness are independent
of $Y$, so its pushforward factors as $\nu_{A'}\otimes P_Y$.
This proves \eqref{eq:factor-density-propagation}. The last assertion
follows from the KL chain rule and the nonnegativity of
$D_{\rm KL}(Q_A\Vert\nu_A)$.
\end{proof}



Suppose initially that a law $Q$ on original weights satisfies
\begin{equation}\label{eq:factor-initial-density}
 Q\le\exp(NK/t^2)\,\mu_{N,t}.
\end{equation}
The virtual prefix records the constructed network, not just its value at
one selected input. Clipping and rounding the first layer are deterministic
channels applied to that original block and initialize the domination
invariant with exponent $NK/t^2$. At an internal stage, let $X$ be the
current original block and $Y$ the remaining original blocks. Induction
with Lemma~\ref{lem:factor-density} gives exponent
\[
 b_s\le NK/t^2+sD_0\le2NK/t^2\qquad(N>n_0).
\]
At that stage abort if its marginal density $u(A,Y)$ is below $\eta$.
Otherwise \eqref{eq:factor-conditional-entropy} and the imported
whole-layer active-direction lemma, applied with tolerance $e$, give
\[
 \dim U\le\frac{t^2}{Ne}
             [b_s+\log(1/\eta)]
       \le\frac{2K+\log(1/\eta)}e\le q.
\]
The active query list is selected from the conditional law given $(A,Y)$,
so it is a function of $(A,Y)$. Its basis is a function of the prefix and
the list, using fixed Gram--Schmidt conventions. The selected list and an
abort symbol admit a fixed alphabet of log-cardinality at most $D_0$.
For reference messages never chosen by the actual selection rule, extend
the basis and transition kernel by fixed measurable default rules. This
makes the uniform-message reference construction in the lemma defined
on its entire domain.

All abort paths remain in the canonical probability law. On an abort,
use a fixed default prefix, retain an unsuccessful flag, and use default
transitions thereafter. Do not condition any intermediate law on success.
The density-test abort is included in the finite message; later aborts
arising from fresh randomness are part of the prefix channel. Consequently
the invariant applies on the full law, and the unsuccessful paths are
discarded only in the final approximation probability estimate.

The domination invariant also provides the old averaged entropy bound.
Thus the imported Gaussian entropy-moment estimate still gives
\begin{equation}\label{eq:factor-active-moment}
 \E\left[\norm{H_0}_F^2/N;\ \text{active success}\right]
       \le4t^2b_s/N+2t^2q\log2\le V.
\end{equation}
In particular, the definitions of $V,B_H,B_1$, and $K_0$ need no extra
factor of $1/\eta$. The improvement in $q$ uses the density hypothesis
\eqref{eq:factor-initial-density}; an arbitrary KL bound alone would not
justify it.

For comparison, using this new rank cutoff with the dense sequential code
of Appendix~\ref{app:all-scales-quantization} already gives
$O(C^6L^{15}F_0^6\log(2+CL))$ samples. The factor construction below saves
additional powers of both resources.

\subsection{A small-rank replacement for the Gaussian field}

We recall the source-sampling notation needed for the new step.
The source activations obey $|h_j(x)|\le1$. If $U$ has $\ell\le q$
orthonormal columns, put $u_j=\sqrt N U_{j,:}$. The half-uniform,
half-leverage law gives weights $a_j=1/(Np_j)\le2$ and
$a_j\norm{u_j}^2\le2\ell$. Draw $m$ source indices $J_b$, and define
\[
 \widehat G=\frac1m\sum_ba_{J_b}u_{J_b}u_{J_b}^{\top},\qquad
 D_{b,:}=\sqrt{a_{J_b}/m}\,u_{J_b}^{\top}\widehat G^{-1/2}.
\]
On the good Gram event, $\widehat G\succeq I/2$ and $D^\top D=I$.
For $\ell=0$, the active matrices are empty and the following projection
is the identity. Write
\[
 P=I-DD^\top,\qquad
 c_b(x)=\sqrt{a_{J_b}/m}\,h_{J_b}(x),\qquad
 \widehat C(x,x)=\norm{Pc(x)}^2\le2.
\]
The ideal field in the imported whitened-sampling estimate has active mean
$H_0\psi(x)$ and independent Gaussian rows of variance $t^2C(x,x)$.
The original reconstruction has mean $H_0\widehat v(x)$ and variance
$t^2\widehat C(x,x)$.

\begin{lemma}[Factored Gaussian reconstruction]
\label{lem:factor-gaussian}
Conditional on the source samples and retained active data, draw an
independent standard Gaussian matrix $R\in\R^{p\times m}$ and independent
standard Gaussian row vectors $g_i\in\R^p$, for all virtual target rows.
Replace the reconstructed coefficients by
\begin{equation}\label{eq:factor-coefficients}
 \widetilde W_{ib}=
   \frac{a_{J_b}}m H_{0,i}\mathbin\cdot u_{J_b}
   +t\sqrt{\frac{a_{J_b}}{mp}}\,(g_i^\top RP)_b.
\end{equation}
This coefficient law does not depend on the input. At every fixed input
it can be coupled to the ideal field so that the expected squared RMS
evaluation error, with the good Gram indicator, is at most
\begin{equation}\label{eq:factor-reconstruction-error}
 \frac{4\norm{H_0}_F^2/N+32(q+1)}m+\frac8p.
\end{equation}
The coefficient matrix in \eqref{eq:factor-coefficients} has rank at most
$\ell+p$.
\end{lemma}
\begin{proof}
Fix the input and abbreviate $\sigma^2=\widehat C(x,x)$.
Conditional on $R$, the new Gaussian evaluations are independent across
rows with variance $t^2s^2$, where
\[
 s=\norm{RPc(x)}/\sqrt p.
\]
If $\sigma>0$, then $s^2/\sigma^2$ has the law of a chi-square variable
with $p$ degrees of freedom divided by $p$. Rationalizing the square root
therefore gives
\begin{equation}\label{eq:factor-sd-error}
 \E_R(s-\sigma)^2
 \le\sigma^2\E(\chi_p^2/p-1)^2
 =2\sigma^2/p\le4/p.
\end{equation}
For $\sigma=0$ the statement is immediate.

To make the coupling explicit, draw $R$ and independent standard normals
$Z_i$, independent of $R$. Use Gaussian row evaluations $t\sigma Z_i$
for the old reconstruction and $tsZ_i$ for the new one. Disintegrating
$g_i$ conditional on its projection along $RPc(x)$ realizes the latter
evaluation while keeping its full conditional marginal standard Gaussian.
Thus the coefficient law is exactly the law in
\eqref{eq:factor-coefficients}. Disintegration likewise extends the old
Gaussian evaluation to its coefficient field. The degenerate zero-variance
case is handled by independent completion. Only this comparison coupling,
not either field marginal, depends on $x$.

The old whitened reconstruction estimate is
$[2\norm{H_0}_F^2/N+16(q+1)]/m$. Couple through that reconstruction and
apply the squared triangle inequality with
\eqref{eq:factor-sd-error}. This yields
\eqref{eq:factor-reconstruction-error}. Finally the active term factors
through $\R^\ell$ and the Gaussian term through $\R^p$, proving the
rank claim.
\end{proof}

Integrating \eqref{eq:factor-reconstruction-error} with the success
indicator and using \eqref{eq:factor-active-moment}, the first line of
\eqref{eq:factor-source-width}, and $p\ge1000/e$, bounds this error by
\[
 \frac{4V+32(q+1)}m+\frac8p<e/8.
\]
The comparison requires no simultaneous coupling of all inputs.
Conditional independence of rows given $R$ suffices; independence of the
new rows after averaging over $R$ is not asserted or needed.

\subsection{Bounded factors and their sequential rounding}

At each internal stage first clip active rows of norm greater than $B_H$
to zero, exactly as in the imported construction. Equation
\eqref{eq:factor-active-moment} and bounded activations give squared RMS
activation cost at most $4V/B_H^2=e$. Apply
Lemma~\ref{lem:factor-gaussian} to this clipped matrix.

Abort if $\norm R_F^2>pm/\eta$. Since $\E\norm R_F^2=pm$, this event
has conditional probability at most $\eta$. Also replace each $g_i$ by
zero when $\norm{g_i}^2>B_g^2=4p/e$. This row replacement has probability
at most $e/4$ and hence expected squared activation change at most $e$:
two tanh values differ by at most two. All virtual rows are clipped before
any later row selection, including selections depending on future weights.

On success, write the resulting coefficients as $A_iB$, with zero padding
to factor dimension $k=q+p$, where
\[
 A_i=(H_{0,i},g_i),\qquad
 B_{{\rm active},ab}=a_{J_b}u_{J_b,a}/m,
\]
\[
 B_{{\rm Gaussian},ab}
       =t\sqrt{a_{J_b}/(mp)}\,(RP)_{ab}.
\]
Every left factor entry has absolute value at most $A_f$. Every right
factor entry has absolute value at most $B_f$, since
\[
 |a_{J_b}u_{J_b,a}/m|\le2\sqrt{q+1}/m,
\qquad
 t\sqrt{\frac{a_{J_b}}{mp}}|(RP)_{ab}|
 \le\sqrt{\frac2{mp}}\norm R_F\le\sqrt{2/\eta}.
\]

Round both factor arrays to the common scalar mesh \eqref{eq:factor-mesh}
before computing the next virtual prefix. For each bounded incoming state,
\begin{equation}\label{eq:factor-rounding-injection}
 |(A_{i,\rm round}B_{\rm round}-A_iB)h|
 \le mk(A_f+B_f+1)\zeta\le\sqrt e/4.
\end{equation}
Here each product changes by at most
$(A_f+B_f+1)\zeta$, because $\zeta\le1$.
Round first-layer coefficients and reconstructed readout coefficients
directly. Their preactivation errors are at most $(d+1)\zeta$ and
$m\zeta$, respectively, again at most $\sqrt e/4$.
The rounded candidate is an ordinary depth-$L$ tanh network: its internal
matrix is the product of two rounded factor arrays. Factorization adds
no activation layer.

All source samples, $R$, and Gaussian row factors are fresh randomness
conditional on the prefix, query list, and retained active matrix.
Clipping, factor multiplication, and rounding are deterministic subsequent
operations. They therefore form a channel of the precise kind required
by Lemma~\ref{lem:factor-density}. The reference-prefix law need not be
Gaussian or have a density on its own coordinate space. The comparison
coupling does not replace this canonical channel in the density induction.

\subsection{The disturbance budget and the finite class}

Consider first $N>n_0$ and a law satisfying
\eqref{eq:factor-initial-density}. The first layer is clipped at the
physical standardized row bound $B_1$ and then rounded. Its expected
squared RMS activation injection is less than $3e$ by the imported
first-layer estimate and the squared triangle inequality.
At an internal stage the five contributions are active-direction
transport, active-row clipping, reconstruction, Gaussian-row clipping,
and factor rounding. Their expected squared RMS bounds, with success
indicators, are respectively
\[
 2e,\quad e,\quad e/8,\quad e,\quad e/16.
\]
Another squared triangle inequality bounds their sum by
$5(2+1+1/8+1+1/16)e<21e$. In particular the conservative bound $41e$
from Appendix~\ref{app:all-scales-quantization} remains available.

The good Gram probability follows from the second line of
\eqref{eq:factor-source-width}. Use the original readout construction:
truncate at $N\norm v^2\le K_0$, sample source indices with probabilities
proportional to $|v_i|$, and use signed coefficients of magnitude
$\norm v_1/m\le\sqrt{K_0}/m$. The last line of
\eqref{eq:factor-source-width} gives output sampling error at most $r/16$
uniformly over the finite input domain except with probability $\eta$.
The energy truncation also fails with probability at most $\eta$.
Including density-test, Gram, and shared-sketch aborts at internal stages,
the total abort probability is at most $4L\eta$.

For each fixed input, glue the conditional transports as in
Appendix~\ref{app:all-scales-quantization}: retain the full original-weight
marginal $Q$ and the canonical outgoing prefix/future-weight marginal at
each transition. Extend the transition by the previous conditional law of
past original weights. This preserves every earlier computed prefix.
The final candidate output is exactly the original network under the
postactivation and readout injections defined by its rounded virtual
computation. No robustness of an intermediate candidate is assumed.

On the non-aborted paths, Markov's inequality gives
\[
 \Prb\left\{\sum_{s<L}\norm{\eta^s}^2/N>r^2/16,
                         \ \text{no abort}\right\}
 \le\frac{656Le}{r^2}=0.656\alpha.
\]
Readout rounding adds at most $\sqrt e/4\le r/64$. Thus, outside probability
at most $0.656\alpha+4L\eta<\alpha$, the total disturbance satisfies
\[
 \norm\eta_*\le
 \sqrt{r^2/16+(r/16+r/64)^2}
       =\frac{\sqrt{281}}{64}r<r/2.
\]

The construction above does not require $m\le N$. Source sampling is with
replacement, and duplicated indices give duplicated stored neurons. The
Gram lemma also allows this: on its good event $D$ has orthonormal columns,
and $\ell\le q<m$. Hence the only width cutoff needed for this branch is
$N>n_0$, which pays for the accumulated message exponents.

For $N\le n_0$, the original density bound implies a KL bound
$NK/t^2$. The Gaussian entropy-moment inequality gives, for all physical
standardized coordinates together,
\[
 \E_Q\norm{tZ_{\rm all}}^2
       \le4n_0K+2P_{\rm small}\log2.
\]
With probability at least $1-\alpha/4$, every actual coefficient is at
most $B_{\rm small}$ in absolute value. On that event round the original
network directly; otherwise use a zero-output default. The bounded rounded
activations give the exact injection bound
\[
 \norm\eta_*^2
 \le\zeta^2[(d+1)^2+(L-1)N^2]
 \le\zeta^2L(n_0+d+1)^2\le Le/16<r^2/4.
\]
Pad these small networks only to width $n_0$.

The large branch keeps the $m$ selected rows at each layer. It stores
$m(d+1)$ first-layer scalars, $2mk$ factor scalars for each internal
matrix, and $m$ readout scalars. The small branch stores at most
$P_{\rm small}$ scalars. Every rounded scalar in either branch belongs
to a fixed alphabet of at most $3+4B_{\rm code}/\zeta$ values.
The union of these two fixed finite classes, denoted $\mathcal H_{\rm fac}$,
therefore satisfies
\begin{equation}\label{eq:factor-finite-class}
 \log|\mathcal H_{\rm fac}|\le J_{\rm fac}.
\end{equation}
In particular one does not count the small branch as a dense network of
width $m$.

We have proved that every law satisfying
\eqref{eq:factor-initial-density} admits a single candidate law
$\Gamma_Q$ on $\mathcal H_{\rm fac}$ with the following property: at each
fixed input it has a coupling with $w\sim Q$ and a disturbance such that
\begin{equation}\label{eq:factor-final-coupling}
 h(x)=z_w(x;\eta),\qquad \norm\eta_*<r/2
\end{equation}
except with probability at most $\alpha$. Both marginals are independent
of the input used to form the comparison coupling.

\subsection{Posterior comparison and polynomial accounting}

\begin{proof}[Proof of Theorem~\ref{thm:factor-all-scales}]
The product target cylinder gives raw conditioning probability at least
$\exp(-ND_*/t^2)$. The chi-square exponential moment for the readout,
combined with this denominator, gives
\[
 Q_{S,N,t,c}\{E_v>B_v\}
 \le\exp\left(-\frac{NB_v}{4t^2}
                  +\frac N2\log2+\frac{ND_*}{t^2}\right)
 \le\beta.
\]
Let $Q'$ be the posterior conditioned on $E_v\le B_v$. Unlike a generic
law controlled only in relative entropy, it satisfies the literal
density bound
\[
 \frac{dQ'}{d\mu_{N,t}}
 \le\frac{\exp(ND_*/t^2)}{1-\beta}
 \le\exp(NK/t^2).
\]
Every original network in its support is robust on the common radius $r$.
Consequently \eqref{eq:factor-final-coupling} implies clipped-score error
at most $1/16$ on its good event.

Use the ramp $\psi$ from the earlier statistical proof: it is one on
$(-\infty,1/3]$, zero on $[5/12,\infty)$, and affine in between.
Every training point has candidate ramp loss zero on the good coupling;
every original test-margin failure has candidate loss one there.
A relative Chernoff bound and a union bound over
$\mathcal H_{\rm fac}$ give, on one sample event of probability at least
$1-\delta$,
\[
 \E_D\psi(f(X)T(h(X)))
 \le\frac2n\sum_{i=1}^n\psi(f(X_i)T(h(X_i)))
       +\frac{8[J_{\rm fac}+\log(1/\delta)]}{n}
 \quad(h\in\mathcal H_{\rm fac}).
\]
Averaging this inequality under $\Gamma_{Q'}$, applying the two marginal
comparisons, and restoring the discarded readout tail bounds the desired
risk by
\[
 \beta+3\alpha+
       \frac{8[J_{\rm fac}+\log(1/\delta)]}{n}\le\varepsilon.
\]
The finite class and sample event do not depend on $N,t$, which gives the
simultaneous conclusion. Additional global clean-margin conditioning retains
the same target cylinder, density bound, and readout-tail argument, so the
same proof applies to that stated optional variant.
\end{proof}

\begin{proof}[Proof of Corollary~\ref{cor:factor-powers}]
The product-cylinder constants give, at fixed ancillary resources,
\[
 K,B_v=O(CLF_0),\qquad r^{-2}=O(CL^2F_0),\qquad
 e^{-1}=O(CL^3F_0),\qquad \eta^{-1}=O(L).
\]
The term $\log(1/\eta)$ in the rank bound is dominated by the target cost
$K$. Therefore
\[
 q,V=O(C^2L^4F_0^2),\qquad p=O(CL^3F_0),\qquad
 m=O(C^3L^7F_0^3),\qquad n_0=O(CL^4F_0).
\]
The Gram and readout terms in the maximum defining $m$, including their
logarithms, are dominated by its first term at this fixed-resource scale.
It follows that
\[
 P_{\rm fac}=O(C^5L^{12}F_0^5),\qquad
 P_{\rm small}=O(C^2L^9F_0^2).
\]
All coefficient bounds and inverse meshes are polynomial in $C,L,F_0$,
so $\log(3+4B_{\rm code}/\zeta)=O(\log(2+CL))$.
Substitution into \eqref{eq:factor-code-budget} and
\eqref{eq:factor-sample-budget} proves the claimed order and both special
cases.
\end{proof}

\begin{remark}[A conservative joint polynomial for the factored code]
\label{rem:factor-joint-polynomial}
Put
\[
 S=2+d+L+C+c+c^{-1}+\varepsilon^{-1}+\log(1/\delta).
\]
The target-cylinder estimates and
\eqref{eq:factor-posterior-budgets} give
\[
 K,B_v=O(S^3),\qquad r^{-2}=O(S^6),\qquad
 e^{-1}=O(S^8),\qquad \eta^{-1}=O(S^2).
\]
The new budgets consequently satisfy
\[
 q,V=O(S^{11}),\quad D_0=O(S^{12}),\quad
 m=O(S^{19}),\quad n_0=O(S^{13}),\quad p=O(S^8),
\]
\[
 P_{\rm fac}=O(S^{31}),\quad P_{\rm small}=O(S^{27}),\quad
 B_{\rm small}^2=O(S^{28}),\quad B_{\rm code}=O(S^{14}).
\]
For the mesh, $A_f=O(S^{10})$ and $B_f=O(S)$; the latter uses the
first lower bound on $m$ to bound $\sqrt{q+1}/m$ by a constant.
Thus $\zeta^{-1}=O(S^{44})$, where the factor-rounding denominator
dominates the other displayed choices. It follows that
\[
 J_{\rm fac}=O(S^{31}\log S),\qquad
 n=O(S^{32}\log S)\subseteq O(S^{33}).
\]
This degree-33 envelope is deliberately conservative. Like the refined
width and depth powers above, it is a consequence of the present
conventional argument and carries no Lean-certification claim.
\end{remark}

\section{Entropy-weighted arithmetic compression}
\label{app:entropy-metric}

This appendix constructs a finite comparison class while leaving the
original Gaussian prior and the robustness support unchanged. The
comparison functions use a finite arithmetic decoder between tanh layers;
they need not be ordinary neural networks. This is sufficient for the
finite-class statistical argument and for the layerwise internal-noise
coupling. The result is a conventional proof, not a Lean formalization.

\subsection{A single entropy matrix for every Gaussian channel}

\begin{lemma}[Gaussian-channel entropy matrix]
\label{lem:em-channel}
Let $P$ be a probability law on $\mathbb R^{N\times q}$ with finite
$D=D_{\rm KL}(P\Vert\gamma_{Nq})$. There is a positive semidefinite
$q\times q$ matrix $\mathsf T$ with $\operatorname{tr}\mathsf T=D$
such that, for
\[
 Q_b=\operatorname{Law}(Xb+\sqrt{1-\|b\|^2}\,Z),
 \qquad X\sim P,\quad Z\sim\gamma_N\text{ independent},
\]
every $\|b\|<1$ and $\|b'\|\le1$ satisfy
\[
 W_2^2(Q_b,Q_{b'})\le
 \frac{8(b'-b)^\top\mathsf T(b'-b)}{1-\|b\|^2}.
 \tag{EM1}
\]
The same matrix works for all channel directions.
\end{lemma}
\begin{proof}
Let $P_s$ be the Ornstein--Uhlenbeck evolution of $P$, let
$f_s=dP_s/d\gamma_{Nq}$, and view its relative score as an
$N\times q$ matrix. Define
\[
 \mathsf I_s=\mathbb E_{P_s}
       [(\nabla\log f_s)^\top(\nabla\log f_s)],
 \qquad \mathsf T=\int_0^\infty\mathsf I_s\,ds.
\]
Gaussian entropy dissipation gives $\operatorname{tr}\mathsf T=D$.
Score commutation and conditional Jensen give, in Loewner order,
\[
 \mathsf I_t\preceq e^{-2(t-s)}\mathsf I_s\quad(0<s<t),
 \qquad
 \mathsf I_t\preceq\frac{2}{e^{2t}-1}\mathsf T.
 \tag{EM2}
\]
First assume a smooth positive density. Put $C=1-\|b\|^2$.
For a direction velocity $\delta$, differentiating
$Y=Xb+\sqrt C Z$ yields the continuity-equation velocity
\[
 v_b(y)=\mathbb E\left[
       X\delta-\frac{b\cdot\delta}{\sqrt C}Z\mid Y=y\right]
       =\mathbb E[\nabla\log f(X)\mid Y=y]\,\delta.
 \tag{EM3}
\]
To verify the equality, Gaussian integration by parts under the
reference conditional law of $X$ given $Y$ gives
\[
 \mathbb E[X\mid Y]-Yb^\top
  =\mathbb E[\nabla\log f(X)\mid Y](I-bb^\top).
\]
Substitution, using $(I-bb^\top)^{-1}=I+bb^\top/C$, proves (EM3).

For $\rho=\|b\|^2>0$, take $t$ such that $e^{-2t}=\rho$.
The channel is a projection of the full observation
$X_t=\sqrt\rho X+\sqrt{1-\rho}G$.
Score commutation and (EM2) imply
\[
 \mathbb E\|v_b(Y)\|^2
 \le\rho^{-1}\delta^\top\mathsf I_t\delta
 \le\frac2C\delta^\top\mathsf T\delta.
 \tag{EM4}
\]
The case $b=0$ follows by continuity. Along the straight path
$b_s=(1-s)b+sb'$, with $\delta=b'-b$, the residual variance is
\[
 C_s=(1-s)C+s(1-\|b'\|^2)+s(1-s)\|\delta\|^2
       \ge(1-s)C.
\]
The Wasserstein length bound and (EM4) give
\[
 W_2(Q_b,Q_{b'})
 \le\sqrt{2\delta^\top\mathsf T\delta}
                   \int_0^1 C_s^{-1/2}\,ds
 \le2\sqrt{2\delta^\top\mathsf T\delta/C}.
\]
For general finite-entropy $P$, apply the result first to $P_\epsilon$.
Its entropy matrix is $\int_\epsilon^\infty\mathsf I_sds\preceq\mathsf T$.
Finite Gaussian relative entropy implies a finite second moment;
OU smoothing and the channel outputs converge in $W_2$ as
$\epsilon\downarrow0$. This proves the stated generality and also
the boundary endpoint $\|b'\|=1$.
\end{proof}

For a physical active matrix $H=tX$, use
\[
 T=\frac{t^2}{N}\mathsf T.
\]
If $\|b\|^2+C=V$, $C>0$, and $\|b'\|^2\le V$, rescaling (EM1)
gives the following squared RMS bound for the two physical channels:
\[
 \frac1N W_2^2\left(
 \operatorname{Law}(Hb+t\sqrt C Z),
 \operatorname{Law}(Hb'+t\sqrt{V-\|b'\|^2}Z')\right)
 \le\frac{8V}{C}(b'-b)^\top T(b'-b).
 \tag{EM5}
\]
In our construction $V\le2$.

\subsection{Budgets and the finite message}

Fix $K\ge1$, $0<\alpha\le1/8$, $0<r\le1/6$, $L\ge2$,
and a scale ceiling $0<t_0\le1$. Set
\[
 \begin{gathered}
 e=\frac{\alpha r^2}{1000L},\qquad
 \eta=\frac{\alpha}{100L},\\
 K_a=2K+\log(1/\eta),\qquad q=\lceil K_a/e\rceil,
 \qquad B=8K_a.
 \end{gathered}
 \tag{EM6}
\]
At a successful internal stage the bounded-density argument and the
whole-layer active-direction lemma give rank at most $q$ and a
conditional physical entropy matrix with $\operatorname{tr}T\le K_a$.
Round its symmetric entries to the nearest multiple of $e/q$ and add
$2eI$. The resulting matrix $\overline T$ satisfies
\[
 \overline T\succeq T+eI,\qquad
 \operatorname{tr}\overline T\le B,\qquad
 \kappa(\overline T)\le B/e.
 \tag{EM7}
\]
Indeed, the rounding error has operator norm at most $e/2$, and
$qe\le K_a+e\le2K_a$. Smaller active ranks are padded by zeros before
this operation, with a separate diagonal floor on the unused block.

Let
\[
 \begin{gathered}
 A_T=3+8K_aq/e,\qquad M=\max\{1,2LD_0/K\},\\
 D_0=[(d+1)q+2]\log2+q^2\log A_T.
 \end{gathered}
 \tag{EM8}
\]
The query list, abort flag, and quantized matrix form a fixed finite
message of log-cardinality at most $D_0$. Invalid reference matrix
messages use a fixed positive definite default and a fixed transition.

For a law
\[
 Q\le\exp(NK/t^2)\mu_{N,t},\qquad 0<t\le t_0,
 \tag{EM9}
\]
the future-density invariant has exponent
$b_s\le NK/t^2+sD_0$. In the branch $N/t^2>M$, this is at most
$2NK/t^2$. The density test discards probability at most $\eta$ at
each stage and bounds conditional current-block entropy by
$2NK/t^2+\log(1/\eta)$. This justifies both $q$ and (EM7).
The matrix message depends on the prefix and future original weights;
after it is supplied the reconstruction kernel has no additional
dependence on those future weights.

\subsection{Independent-batch regression in the entropy metric}

At a stage, let $\ell\le\min(q,N)$ be its actual active rank. All
regression and projection formulas in this subsection and the next
operate on the $\ell$-dimensional active block; $q$ remains only the
budget, mesh, and counting cap. In these formulas $\overline T$ denotes
the corresponding active block. Coordinates are padded only for final
storage, with a separate diagonal metric on the unused block, so that
the padded projection keeps all unused coordinates zero.
Write $u_j=\sqrt N U_{j,:}$ for the active source features and, at a
fixed input, put
\[
 \psi=\frac1N\sum_j u_jh_j,\quad
 K_x=\frac1N\sum_jh_j^2,\quad
 R_j=h_j-u_j^\top\psi,\quad C=\frac1N\sum_jR_j^2.
\]
Then $N^{-1}\sum_j u_ju_j^\top=I_\ell$,
$N^{-1}\sum_j u_jR_j=0$, and $K_x=\|\psi\|^2+C\le1$.

Choose
\[
 m_0=\left\lceil\max\{16q,
       64(q+1)\log(4(q+1)/\eta)\}\right\rceil,
 \qquad m_A=\lceil10^5B/e\rceil.
 \tag{EM10}
\]
Use $m_0$ ordinary half-uniform, half-leverage samples for a pilot
$\psi_0=G_0^{-1}v_0$. Abort if $G_0\not\succeq I/2$.
The usual matrix Chernoff bound makes this probability at most $\eta$.
With the good-pilot indicator,
\[
 \mathbb E\|\psi_0-\psi\|^2\le8qC/m_0\le C/2.
 \tag{EM11}
\]

Independently draw $m_A$ samples with probabilities
\[
 p_j=\frac1{3N}\left(1+\frac{\|u_j\|^2}{\ell}
       +\frac{u_j^\top\overline T u_j}
                    {\operatorname{tr}\overline T}\right),
 \qquad a_j=(Np_j)^{-1}.
 \tag{EM12}
\]
Empty active ranks use uniform sampling and empty matrices. In positive
rank, (EM12) gives
\[
 a_j\le3,\qquad a_j\|u_j\|^2\le3q,\qquad
 a_j u_j^\top\overline T u_j\le3\operatorname{tr}\overline T.
\]
Define
\[
 \widehat\psi=\psi_0+\frac1{m_A}
       \sum_b a_{J_b}u_{J_b}(h_{J_b}-u_{J_b}^\top\psi_0).
 \tag{EM13}
\]
Conditional on the pilot this estimator is unbiased for $\psi$.
Since
$N^{-1}\sum_j(h_j-u_j^\top\psi_0)^2=C+\|\psi_0-\psi\|^2$,
\[
 \mathbb E[\|\widehat\psi-\psi\|_{\overline T}^2;
                  \text{good pilot}]
 \le\frac{3BC}{m_A}\left(1+\frac{8q}{m_0}\right)
 \le\frac{9BC}{2m_A}.
 \tag{EM14}
\]
This calculation does not assume the false isotropic residual bound
$\mathbb E bb^\top\preceq2CI/m$. The independent correction is
what avoids conjugating the entropy metric by a random Gram inverse.
When $C=0$, the pilot and corrected estimator are both exactly $\psi$
on the good pilot event.

The estimator is linear in its sampled activation values. Write it as
$\widehat\psi(x)=Fh_{\rm src}(x)$, with zero columns for unused source
groups. A deterministic entry bound on success is
\[
 |F_{ab}|\le F_*:=20(q+1)^2.
 \tag{EM15}
\]
For example, pilot columns have norm at most $4\sqrt q/m_0$;
the second empirical Gram has norm at most $3q$; and correction columns
have norm at most $3\sqrt q/m_A$. These facts imply (EM15).

\subsection{A positive variance envelope and the canonical row law}

Set
\[
 \lambda=\min\{1,e/(10^5t_0^2)\},\qquad
 m_K=\max\left\{1,\left\lceil
                \frac{\log(2/\eta)}{2\lambda^2}\right\rceil\right\}.
 \tag{EM16}
\]
Take a third independent uniform source batch and define
\[
 \widehat K(x)=\frac1{m_K}\sum_b h_{J_b}(x)^2,
 \qquad K_+(x)=\widehat K(x)+\lambda.
\]
At each fixed input, Hoeffding gives
\[
 \Pr\{|\widehat K-K_x|>\lambda\}\le\eta.
 \tag{EM17}
\]
For $\lambda=1$ the event is empty. On its complement,
$K_x\le K_+\le K_x+2\lambda$, and always $0<K_+\le2$.
This input-specific event is discarded only in the comparison; the
canonical function is defined at every input.

Define the metric projection
\[
 a(x)=\mathop{\rm argmin}_{\|z\|^2\le K_+(x)}
                   \|z-\widehat\psi(x)\|_{\overline T}^2,
 \qquad \sigma(x)=\sqrt{K_+(x)-\|a(x)\|^2}.
 \tag{EM18}
\]
The minimizer is unique. On the event in (EM17), $\psi$ is feasible;
nonexpansiveness of Hilbert-space projection gives
\[
 \|a-\psi\|_{\overline T}^2
       \le\|\widehat\psi-\psi\|_{\overline T}^2.
 \tag{EM19}
\]
The unrounded virtual target row is
\[
 h_i^{\rm new}(x)=\tanh\{H_i a(x)+t g_i\sigma(x)\},
 \qquad g_i\sim N(0,1)\text{ independently over rows}.
 \tag{EM20}
\]
One scalar Gaussian per row suffices. Its cross-input covariance need
not match the original field: every comparison is at one fixed input,
where (EM20) has the required independent Gaussian row marginals.
The complete function law in (EM20) is nevertheless input-independent.

To bound its error, put $\Delta=K_+-K_x\in[0,2\lambda]$ on (EM17).
Increasing the ideal residual variance from $C$ to $C+\Delta$ costs
at most $t^2\Delta\le2t_0^2\lambda$ in squared RMS. The enlarged ideal
field and (EM20) have common Gaussian-reference variance $K_+\le2$.
For $C>0$, (EM5), (EM19), and (EM14) therefore give, after the squared
triangle inequality and averaging,
\[
 \mathbb E[\text{squared RMS reconstruction error};\text{good}]
 \le4t_0^2\lambda+\frac{144B}{m_A}<e/8.
 \tag{EM21}
\]
The inequality uses $T\preceq\overline T$ and $C+\Delta\ge C$.
For $C=0$, exact interpolation makes $a=\psi$, and only the added
variance remains. Thus (EM21) also covers that case.

The channel coupling need not preserve the realized $H$. Conditional
on prefix and future weights, both sides have the same specified
active-matrix marginal. Disintegration extends the evaluation coupling
to the original current block on one side and to the complete canonical
candidate field on the other. This preserves the outgoing
prefix/future marginal required by the sequential gluing argument.

\subsection{Clipping, an explicit mesh, and the decoder size}

Retain the old first-layer and readout budgets, and set
\[
 V_H=8K+2q,\quad B_H^2=4V_H/e,\quad B_Z=2/\sqrt e,
 \quad B_1^2=4[4K+2(d+1)\log2]/e,
\]
\[
 \begin{gathered}
 K_0=(4K+2\log2)/\eta,\qquad m=m_0+m_A+m_K+m_R,\\
 m_R=\left\lceil512K_0r^{-2}
       [d\log2+\log(2/\eta)+1]\right\rceil.
 \end{gathered}
 \tag{EM22}
\]
Clip an active row to zero when its norm exceeds $B_H$, and clip the
physical scalar $tg_i$ to zero when its absolute value exceeds $B_Z$.
The old entropy-moment estimate makes the active clipping cost at most
$e$ in expected squared activation RMS. Since
$\Pr\{|tg_i|>B_Z\}\le e/4$, Gaussian clipping has the same bound.
Both clippings occur for all virtual rows before any future row selection.
The entropy matrix is computed from the unclipped active law; clipping
is paid after the comparison in (EM21).

Put
\[
 A_* =\sqrt{B/e}\,m\sqrt{q+1},\qquad
 U_*=2\sqrt{q+1}+B_HA_*+2,
\]
\[
 \zeta=\min\left\{1,\frac{\sqrt e}{8U_*},
       \frac{e}{256B_Z^2 A_*},
       \frac{\sqrt e}{4(d+1)},\frac{\sqrt e}{4m}\right\}.
 \tag{EM23}
\]
Round the retained active scalars, entries of $F$, and physical Gaussian
scalars to this common mesh. The matrix $\overline T$ stays on its
defining finite grid. The current incoming activations are already those
of the rounded canonical prefix, so $\widehat K$ and the projection
radius are unchanged when performing this local rounding comparison.

Rounding $F$ changes its feature vector by at most
$m\sqrt{q+1}\zeta$. Projection is $\sqrt{B/e}$-Lipschitz in Euclidean
norm, so its output changes by at most $A_*\zeta$. Both projected
vectors have norm at most $\sqrt2$, and hence
\[
 |\sigma(a)-\sigma(a')|^2
 \le|\|a\|^2-\|a'\|^2|\le3\|a-a'\|.
\]
For every bounded incoming state, the complete preactivation rounding
error is therefore at most
\[
 U_*\zeta+2B_Z\sqrt{A_*\zeta}\le\sqrt e/4.
 \tag{EM24}
\]
First-layer and readout rounding have the same bound from the last two
terms of (EM23). All meshes and coefficient bounds are polynomial in
the theorem resources.

The decoder computes $Fh$, the average of the designated source squares,
the metric projection (EM18), the scalar square root, and then (EM20).
The projection is an exact operation of this finite arithmetic decoder;
no computational-complexity assertion is needed. Let
\[
 B_{\rm code}=1+\max\{B_1,B_H,B_Z,F_*,\sqrt{K_0},1\},
\]
\[
 P=m(d+1)+(L-2)(2mq+m)+m,
\]
\[
 J_{\rm EM}=\log2+P\log(3+4B_{\rm code}/\zeta)
                  +(L-2)q^2\log A_T.
 \tag{EM25}
\]
Padding unused source groups and rows gives a single finite function
class $\mathcal H_{\rm EM}$ with
$\log|\mathcal H_{\rm EM}|\le J_{\rm EM}$, independent of $N$ and of
the actual scale $0<t\le t_0$.

\subsection{Compression and resource consequences}

\begin{theorem}[Arithmetic compression]
\label{thm:em-compression}
With (EM6)--(EM25), suppose (EM9) and $N/t^2>M$. There is one law
$\Gamma_Q$ on $\mathcal H_{\rm EM}$ such that, at every fixed Boolean
input, it can be coupled to the original law $Q$ and an internal
disturbance $\xi$ satisfying
\[
 h(x)=z_w(x;\xi),\qquad \|\xi\|_*<r/2
\]
except with probability at most $\alpha$. The two specified marginals
do not depend on the input used for the comparison.
\end{theorem}
\begin{proof}
The active-direction transport, active clipping, reconstruction,
Gaussian clipping, and rounding contributions have expected squared
RMS at most $2e,e,e/8,e,e/16$, respectively. The squared triangle
inequality retains the old conservative $41e$ bound per hidden stage.
The first-layer bound is less than $3e$. Density tests and pilot Gram
failures are canonical aborts; variance-envelope failures are
input-specific comparison failures. Include every abort path in the
canonical law and use default transitions thereafter. Along with the
old readout truncation and sampling failures their total probability is
at most $4L\eta$.

Every new-prefix operation uses only the prefix, finite message,
retained original active matrix, and fresh randomness. Thus the
bounded-density induction is valid on the full law. The conditional
evaluation couplings glue while preserving the original-weight marginal
and each outgoing canonical prefix/future marginal. The original
network evaluated under the resulting postactivation injections has
exactly the candidate output. No robustness of a virtual prefix is used.

The old disturbance calculation now applies unchanged:
\[
 \Pr\left\{\sum_{s<L}\|\xi^s\|^2/N>r^2/16,
                         \ \text{no comparison failure}\right\}
 \le656Le/r^2=0.656\alpha.
\]
The readout approximation and rounding contribute at most
$r/16+r/64$. Consequently the total disturbance is less than $r/2$
outside probability $0.656\alpha+4L\eta<\alpha$.
\end{proof}

The complementary branch $N/t^2\le M$ should be handled by the raw-prior
posterior generalization bound, whose relative-entropy charge is at most
$NK/t^2\le MK=O(LD_0)$. A dense parameter code in that branch would
lose the improvement. Uniformity in width and scale is supplied by the
separate finite Gaussian-scale domination argument. Since $m\ge q$,
this charge and (EM25) are both bounded by the coarser expression
$\widetilde O(L[q^2+(d+q)m])$.

Ignoring only logarithmic factors and fixed ancillary resources,
\[
 q=\widetilde O(K/e),\qquad
 m=\widetilde O\left(K/e+t_0^4/e^2+m_R\right),
\]
\[
 J_{\rm EM}+LD_0
 =\widetilde O\left(LK^2/e^2+t_0^4LK/e^3\right).
 \tag{EM26}
\]
Here the displayed orders also assume the original readout count is
dominated by $\widetilde O(K/e)$, as it is for the posterior budgets
under consideration. At $t_0=1$, the product-cylinder bounds
$K=O(CLF_0)$ and $e^{-1}=O(CL^3F_0)$ give
\[
 J_{\rm EM}+LD_0=\widetilde O(C^4L^{11}F_0^4).
\]
The original ramp-loss comparison applies to this finite function class,
using the same posterior readout truncation and common robustness radius.

\section{Direct posterior bounds and interpolation between Gaussian scales}
\label{app:scale-interpolation}

The finite message used by the entropy-metric reconstruction has a
quadratic dependence on the active rank. Its accumulated cost need not
be converted into a dense small-width code. The actual conditioning
posterior can be treated directly whenever that message cost is larger
than its Gaussian relative-entropy bound. We give a single statistical
event for this argument and for the scale-dependent reconstruction.
All reference laws below are fixed before the training sample is drawn.
This section prices the comparison class uniformly. The later
Gaussian-weighted construction improves its all-scale rate and removes
the intermediate exponent restriction appearing in this preliminary
calculation.

\begin{lemma}[A relative-entropy risk inequality]
\label{lem:si-risk}
Let $\Pi$ be a probability law on a measurable family of predictors,
and let $\ell_h(x)\in[0,1]$. Write
$R(h)=\mathbb E_X\ell_h(X)$ and
$\widehat R_S(h)=n^{-1}\sum_{i=1}^n\ell_h(X_i)$.
With sample probability at least $1-\delta$, simultaneously for every
probability law $\Gamma\ll\Pi$,
\begin{equation}\label{eq:si-risk}
 \mathbb E_\Gamma R
 \le 2\mathbb E_\Gamma\widehat R_S
       +\frac{2}{n}
          \left[D_{\rm KL}(\Gamma\Vert\Pi)+\log(1/\delta)\right].
\end{equation}
\end{lemma}
\begin{proof}
Put $a=1-\exp(-1)>1/2$. Convexity gives
$e^{-u}\le1-au$ for $0\le u\le1$. For each fixed predictor,
independence of the sample therefore gives
\[
 \mathbb E_S\exp\{n[aR(h)-\widehat R_S(h)]\}
 \le e^{naR(h)}[1-aR(h)]^n\le1.
\]
Fubini and Markov's inequality imply, with probability at least
$1-\delta$,
\[
 \int\exp\{n[aR(h)-\widehat R_S(h)]\}\,d\Pi(h)
 \le1/\delta.
\]
On this event, Jensen's inequality applied under $\Gamma$ gives
\[
 n\left[a\mathbb E_\Gamma R-
                   \mathbb E_\Gamma\widehat R_S\right]
 \le D_{\rm KL}(\Gamma\Vert\Pi)+\log(1/\delta).
\]
Indeed, the general inequality used here is
$\int F\,d\Gamma-D_{\rm KL}(\Gamma\Vert\Pi)
 \le\log\int e^F\,d\Pi$; it follows by applying concavity of the
logarithm to $e^F(d\Pi/d\Gamma)$. Division by $a>1/2$ proves
\eqref{eq:si-risk}. Infinite relative entropy makes the assertion
vacuous and causes no exception.
\end{proof}

The target cylinder gives, for the actual posterior,
\begin{equation}\label{eq:si-raw-density}
 \frac{dQ_{S,N,t,c}}{d\mu_{N,t}}
 \le\exp(NK/t^2).
\end{equation}
Here $K\ge1$ may be the same sample-independent target-cost bound used
by the reconstruction. In particular,
$D_{\rm KL}(Q_{S,N,t,c}\Vert\mu_{N,t})\le NK/t^2$.
For the loss
\[
 \ell_w(x)=\mathbf1\{f(x)z_w(x)\le1/4\},
\]
the empirical loss is zero under this posterior. The following finite
reference family makes the resulting risk bound simultaneous in width
and scale; an infinite mixture over widths or exact scales is unnecessary.

\begin{lemma}[A finite Gaussian reference family]
\label{lem:si-grid}
Fix $K\ge1$ and $J\ge K$. Set
\[
 M=\lceil J/K\rceil,\qquad
 P_M=M(d+1)+(L-2)M^2+M,
\]
\[
 R_\tau=1+\left\lfloor\frac{P_M}{2}\log(J/K)\right\rfloor,
 \qquad \tau_j=\exp(-j/P_M)
          \quad(0\le j<R_\tau).
\]
There is a sample-independent prior on the disjoint union of these
$MR_\tau$ Gaussian models such that, whenever $NK/t^2\le J$, the
original posterior, embedded in one component, has relative entropy at
most
\begin{equation}\label{eq:si-grid-kl}
 NK/t^2+1+\log(MR_\tau).
\end{equation}
\end{lemma}
\begin{proof}
The hypothesis implies $N\le M$ and $t^2\ge K/J$.
Choose $j=\lfloor P_M\log(1/t)\rfloor$. Then
$0\le j<R_\tau$, $\tau_j\ge t$, and
$\log(\tau_j/t)\le1/P_M$.
At a fixed width all Gaussian covariance matrices are multiplied by
the same scalar $t^2$. Their density ratio is therefore
\[
 \log\frac{d\mu_{N,t}}{d\mu_{N,\tau_j}}(w)
 =P_N\log(\tau_j/t)
   -\frac N2(t^{-2}-\tau_j^{-2})\kappa_N(w)
 \le P_N\log(\tau_j/t)\le1.
\]
The nonnegative weighted energy $\kappa_N$ is the one defined in the
main text. Combining this pointwise inequality with
\eqref{eq:si-raw-density} bounds relative entropy with respect to
$\mu_{N,\tau_j}$ by $NK/t^2+1$.
Choose the component $(N,j)$ uniformly from the fixed finite family,
then draw from its Gaussian law. Embedding the posterior in the chosen
component adds exactly $\log(MR_\tau)$ and proves the claim.
\end{proof}

We next state the statistical interface needed from the reconstruction.
Let $\psi$ be the ramp loss used in the posterior comparison, and write
\[
 R_\psi(\Gamma)=\mathbb E_{h\sim\Gamma,X}
                         \psi(f(X)T(h(X))),\qquad
 \widehat R_{\psi,S}(\Gamma)=\frac1n\sum_{i=1}^n
                 \mathbb E_{h\sim\Gamma}\psi(f(X_i)T(h(X_i))).
\]
The readout-retention and pointwise coupling estimates give the two
inequalities in the next theorem. Thus its hypotheses can be checked
without imposing robustness on any decoded predictor.

\begin{theorem}[A simultaneous width and scale bound]
\label{thm:si-envelope}
Suppose \eqref{eq:si-raw-density} holds. Fix nonnegative $\alpha,\beta$
and a message threshold $\Delta$. Suppose there are sample-independent
constants
\[
 A\ge\max\{1,K,\Delta\},\qquad B\ge0,
\]
with the following property. For every fixed $0<a\le1$ there is a finite
class $\mathcal H_a$ satisfying
\begin{equation}\label{eq:si-class}
 \log|\mathcal H_a|\le A+Ba^4.
\end{equation}
Whenever $t\le a$ and $NK/t^2>\Delta$, the reconstruction supplies a law
$\Gamma$ on $\mathcal H_a$ such that
\begin{equation}\label{eq:si-comparison}
 \widehat R_{\psi,S}(\Gamma)\le\alpha,
 \qquad
 R(Q_{S,N,t,c})\le\beta+\alpha+R_\psi(\Gamma).
\end{equation}

Set $J=A+B$, use $M,R_\tau$ from Lemma~\ref{lem:si-grid}, and define
\[
 j_a=\left\lceil\frac12\log_2\max\{1,B/A\}\right\rceil,
 \qquad R_a=j_a+1,
 \qquad
 \Lambda=1+\log\frac{2MR_\tau R_a}{\delta}.
\]
With sample probability at least $1-\delta$, simultaneously for every
$N\ge C$ and $0<t\le1$,
\begin{equation}\label{eq:si-envelope}
 R(Q_{S,N,t,c})
 \le\beta+3\alpha+
   \frac8n\left[\min\{NK/t^2,\ A+Bt^4\}+\Lambda\right].
\end{equation}
On the same event, the branch $NK/t^2\le\Delta$ obeys the sharper direct
bound
\begin{equation}\label{eq:si-small}
 R(Q_{S,N,t,c})\le\frac2n[\Delta+\Lambda].
\end{equation}
\end{theorem}
\begin{proof}
Use the finite scale list $a_j=2^{-j/2}$, $0\le j<R_a$.
For every $0<t\le1$, one of these upper scales satisfies $t\le a_j$ and
\begin{equation}\label{eq:si-scale-grid}
 A+Ba_j^4\le4(A+Bt^4).
\end{equation}
Indeed, above the smallest scale choose the nearest upper scale, whose
square is less than twice $t^2$. Below the smallest scale,
$Ba_{j_a}^4\le A$ if $B>A$, while if $B\le A$ the single scale $a=1$
already costs at most $2A$.

Construct one prior as follows. Give weight $1/2$ to the Gaussian
reference family of Lemma~\ref{lem:si-grid}. Give the remaining weight
$1/2$ to a uniform choice of $j$, followed by the uniform law on
$\mathcal H_{a_j}$. Keep the component tag as part of the hypothesis.
For Gaussian components use the test-margin indicator loss; for finite
code components use the ramp loss. These are a single $[0,1]$-valued loss
on the tagged space. Apply Lemma~\ref{lem:si-risk} to this prior.

Put $b=NK/t^2$ and $F=A+Bt^4\le J$.
If $b\le F$, then $b\le J$, so the Gaussian family applies. The
posterior has zero empirical indicator loss, and its relative-entropy
charge, including the mixture weight and confidence term, is at most
$b+\Lambda$. Thus its risk is at most $2(b+\Lambda)/n$.

If $b>F$, then $b>A\ge\Delta$, so the reconstruction applies with the
scale selected in \eqref{eq:si-scale-grid}. Its relative-entropy charge
is at most $4F+\Lambda$. Lemma~\ref{lem:si-risk} and
\eqref{eq:si-comparison} give risk at most
$\beta+3\alpha+2(4F+\Lambda)/n$.
These two cases imply \eqref{eq:si-envelope}.
Finally, $b\le\Delta$ implies $b\le A\le J$, so the direct Gaussian
calculation gives \eqref{eq:si-small} on this same event.
\end{proof}

In particular, taking $\beta=\varepsilon/4$ and
$\alpha=\varepsilon/12$, a sufficient sample condition at any given
width and scale is
\begin{equation}\label{eq:si-sample-envelope}
 n\ge\frac{16}{\varepsilon}
       \left[\min\{NK/t^2,\ A+Bt^4\}+\Lambda\right].
\end{equation}
The event in Theorem~\ref{thm:si-envelope} is common to all such pairs.
For polynomially bounded $A,B,K,d,L$, the additional term $\Lambda$
is logarithmic in these resources and in $1/\delta$.

\begin{corollary}[Interpolation for stronger Gaussian scales]
\label{cor:si-gamma}
Under the hypotheses of Theorem~\ref{thm:si-envelope}, let $\gamma\ge1$.
On its same sample event, it suffices, simultaneously for every $N\ge C$
and $0<t^2\le N^{1-\gamma}$, to require
\begin{equation}\label{eq:si-generic-gamma}
 n\ge\frac{16}{\varepsilon}
 \left[A+
   K^{(2\gamma-2)/(3\gamma-2)}
   B^{\gamma/(3\gamma-2)}+\Lambda\right].
\end{equation}
The convention for $B=0$ is that its displayed contribution is zero.
\end{corollary}
\begin{proof}
If $B\le K$, then $A+Bt^4\le A+B$ and
$B\le K^{(2\gamma-2)/(3\gamma-2)}B^{\gamma/(3\gamma-2)}$;
the result follows directly. Suppose $B>K$ and set
\[
 M_\gamma=(B/K)^{1/(3\gamma-2)},
 \qquad u_0=M_\gamma^{1-\gamma},
 \qquad u=t^2.
\]
If $u\le u_0$, the code term is at most $A+Bu_0^2$.
If $u>u_0$ and $\gamma>1$, the restriction
$u\le N^{1-\gamma}$ implies $N<M_\gamma$, and therefore
$NK/u\le KM_\gamma/u_0=KM_\gamma^\gamma$.
These two terms agree because
\[
 Bu_0^2=KM_\gamma^\gamma
       =K^{(2\gamma-2)/(3\gamma-2)}B^{\gamma/(3\gamma-2)}.
\]
For $\gamma=1$, $u_0=1$ and only the first case is needed.
Thus the minimum in \eqref{eq:si-sample-envelope} is at most the bracket
in \eqref{eq:si-generic-gamma}, proving the claim.
\end{proof}

\paragraph{Application to the entropy-metric code.}
The new message has log-cardinality
\[
 D_0=[(d+1)q+2]\log2+
       q^2\log(3+8K_aq/e),\qquad
 K_a=2K+\log(1/\eta),\quad q=\lceil K_a/e\rceil.
\]
Take $\Delta=\max\{K,2LD_0\}$. The large-branch condition in
Theorem~\ref{thm:si-envelope} is exactly the condition needed to absorb
the accumulated messages under the conservative cutoff
$N/t^2>\max\{1,2LD_0/K\}$.
The complementary branch costs at most $\Delta=O(K+LD_0)$ by
\eqref{eq:si-small}, instead of being rounded as a dense network of
width $LD_0/K$.

For a fixed upper scale $a$, the variance padding can be chosen as
$\lambda=\min\{1,e/(c_0a^2)\}$, where $c_0$ is a sufficiently large
absolute constant. A separate uniform source sample of size
$\widetilde O(1+\lambda^{-2})=
 \widetilde O(1+a^4/e^2)$ estimates the total variance to accuracy
$\lambda$. The extra physical Gaussian variance is $O(a^2\lambda)=O(e)$.
Together with the active width $\widetilde O(K/e)$ and the quadratic
message, the construction therefore gives constants in
\eqref{eq:si-class} with
\begin{equation}\label{eq:si-AB}
 A=\widetilde O(LK^2/e^2),\qquad
 B=\widetilde O(LK/e^3),
\end{equation}
where $A$ is enlarged to include $\Delta$ and the fixed lower-order terms.
These are bounds on pre-sample constants, not on sample-dependent
descriptions. In particular, the direct small branch is absorbed by
the first term in \eqref{eq:si-AB}.

For fixed $d,c,\varepsilon$, put
$F_0=1+\log^2(2+C)$ and $G=\log(2+CL)$.
The target-cylinder estimates give
$K=O(CLF_0)$ and $e^{-1}=O(CL^3F_0)$.
Consequently $q=O(C^2L^4F_0^2)$,
$D_0=O(q^2G)$, and the total source width at upper scale $a$ is
\[
 m=O\left([C^2L^4F_0^2+a^4C^2L^6F_0^2]G\right).
\]
The scalar count is $O(L[mq+q^2])$, up to lower-order first-layer
terms. Polynomial coefficient bounds and inverse meshes contribute one
further factor $G$, uniformly in $a\le1$: in particular
$\lambda^{-1}\le c_0/e$. Thus, after increasing a constant depending
only on $d,c,\varepsilon$, the choices in \eqref{eq:si-AB} satisfy
\begin{equation}\label{eq:si-resources}
 A=O(C^4L^9F_0^4G^2),\qquad
 B=O(C^4L^{11}F_0^4G^2).
\end{equation}
Moreover $\Lambda=O(G+\log(1/\delta))$.
The width and scale envelope becomes
\[
 \widetilde O\min\{CLN/t^2,\ C^4L^9+t^4C^4L^{11}\}.
\]
Corollary~\ref{cor:si-gamma} gives the sufficient sample order
\begin{equation}\label{eq:si-new-powers}
 \widetilde O\left(
 C^4L^9+
 C^{(6\gamma-2)/(3\gamma-2)}
 L^{(13\gamma-2)/(3\gamma-2)}\right).
\end{equation}
An explicit common logarithmic factor in \eqref{eq:si-new-powers} is
$F_0^4G^2$, with an additional additive $O(\log(1/\delta))$ term
before the accuracy-dependent constant. Indeed, the interpolated term
has $F_0$ exponent $(6\gamma-2)/(3\gamma-2)\le4$ and $G$ exponent
$2\gamma/(3\gamma-2)\le2$ for every $\gamma\ge1$.
At $\gamma=1$ this is $\widetilde O(C^4L^{11})$.
For every $\gamma\ge1$ the second term's width exponent is at most four;
its depth exponent is at most nine exactly when $\gamma\ge8/7$.
Thus $\widetilde O(C^4L^9)$ suffices throughout the regime
$0<t^2\le N^{-1/7}$. The existing stronger-scale theorem remains
sharper at $0<t^2\le N^{-1}$, where it gives
$\widetilde O(C^4L^7)$.

The direct posterior argument uses only the target-cylinder denominator
and zero empirical margin loss. It therefore also applies to the optional
global clean-margin conditioning whenever the same cylinder is retained.

\section{A Gaussian-weighted code for uniformly sampled rows}
\label{app:weighted-entropy}

The number of uniformly sampled variance-estimation rows is larger than
that of the pilot and correction rows. Giving every quantized coefficient
the same coding probability is unnecessarily expensive for the uniform
rows. This appendix uses their Gaussian reference probabilities. The
construction still uses the clipped and rounded virtual computation;
there is no change to the prior support or to the prediction convention.
In particular, a uniform sample is taken from a completed virtual layer,
not used to decide which of that layer's rows will be rounded.

\subsection{Uniform subsets and a fixed reference law}

\begin{lemma}[Relative entropy of a uniform row subset]
\label{lem:weighted-shearer}
Let $Y=(Y_1,\ldots,Y_N)$ have a law $P$ on a finite product space, and
let $\rho$ be a probability reference law for one row. If $I$ is a
uniform $k$-element subset of $\{1,\ldots,N\}$ independent of $Y$,
with $1\le k\le N$, list its elements in increasing order. Then
\[
 D_{\rm KL}(\operatorname{Law}(Y_I)\Vert\rho^{\otimes k})
 \le\frac{k}{N}D_{\rm KL}(P\Vert\rho^{\otimes N}).
\]
The same assertion holds conditionally on any variables under which $I$
is independent and uniform. If the finite rows are obtained by applying
the same measurable quantizer to continuous rows, the right side may
be replaced by their relative entropy before quantization.
\end{lemma}
\begin{proof}
For a fixed subset $I$, the chain rule and the fact that extra
conditioning increases average relative entropy give
\[
 D(P_{Y_I}\Vert\rho^{\otimes k})
 \le\sum_{i\in I}\mathbb E
 D(P_{Y_i\mid Y_1,\ldots,Y_{i-1}}\Vert\rho).
\]
Average over the uniform subset. Every index occurs with probability
$k/N$, and the full chain rule gives the asserted right side.
Convexity in the first argument removes the subset index from the law
on the left. Conditional application and data processing prove the
remaining assertions. Infinite relative entropy makes the inequality
vacuous; all uses below have finite relative entropy.
\end{proof}

Use the budgets (EM6)--(EM25) with a fixed upper scale
$0<t_0\le1$, and put
\[
 k=m_K,\qquad m_b=m_0+m_A+m_R,\qquad m=k+m_b.
\]
We require $N\ge k$ in this construction. Each variance batch is now
sampled uniformly \emph{without replacement}. Its average obeys the
same Hoeffding bound as in (EM17): for values in $[0,1]$, the exponential
moment of a sum sampled without replacement is no larger than the
corresponding moment with replacement. Indeed, this comparison is the
inequality between an elementary symmetric mean of the positive numbers
$e^{\theta y_i}$ and their arithmetic mean to the same power.
Thus the existing value of $m_K$ and the existing variance-padding and
reconstruction estimates remain valid.

There is one minor padding convention. At an internal layer of actual
active rank $\ell\le q$, append $q-\ell$ independent Gaussian
coordinates of variance $t^2$ to every active row. The corresponding
rows of $F$ are zero, and the unused block of the metric is diagonal;
the metric projection therefore keeps those coordinates zero. Clip the
\emph{full} $q$-dimensional row at $B_H$, then round all its coordinates
at mesh $\zeta$. The extra coordinates enter no unrounded preactivation.
Their possible effect through clipping is covered by the existing bound
$V_H=8K+2q$. The Gaussian scalar $tg_i$ is clipped and rounded as before.
The resulting row quantizer is fixed: it does not depend on the actual
rank, input, Gaussian scale, learner width, or finite message.

The first-layer raw coordinates are $\sqrt d\,W^1_i$, whose Gaussian
reference covariance is $t^2I_{d+1}$. Its fixed quantizer clips at $B_1$,
rescales by $1/\sqrt d$ as needed for the stored weights, and rounds at
mesh $\zeta$. Let $\Psi_1$ denote this quantizer and let $\Psi_h$ denote
the fixed internal-row quantizer just described. For $u=t^2$ define
\[
 \rho_{1,u}=(\Psi_1)_\#\mathcal N(0,uI_{d+1}),\qquad
 \rho_{h,u}=(\Psi_h)_\#\mathcal N(0,uI_{q+1}).
\]
Both are laws on fixed finite alphabets. The internal quantizer includes
both the padded active vector and its Gaussian scalar.

After constructing a complete virtual layer, independently draw its
uniform $k$-row variance batch. Record the quantized coefficient rows as
$Z_s$, in increasing order of their original indices. If the last hidden
layer does not need a variance batch, draw an unused one there as well;
this harmless convention keeps the dimensions fixed. The reference law
for all these records is
\[
 \Pi_u^Z=\rho_{1,u}^{\otimes k}
          \otimes\bigotimes_{s=2}^{L-1}\rho_{h,u}^{\otimes k}.
\]
It is a fixed pushforward of a centered Gaussian vector with covariance
$uI_p$, where
\begin{equation}\label{eq:weighted-gaussian-dimension}
 p=k(d+1)+(L-2)k(q+1).
\end{equation}
This reference law is independent of the learner width and of all the
other encoded parameters.

\subsection{A joint entropy bound through the canonical construction}

The auxiliary prefix is enriched to retain the coefficient data and
fresh indices from completed stages, so that $Z_1,\ldots,Z_{s-1}$ are
measurable from the incoming prefix $A_s$. One may retain all active
coordinates and Gaussian reconstruction variables that were supplied to
the prefix kernel. This does not mean retaining the discarded original
orthogonal block: an input-specific transport need not preserve its
realization. The retained data are exactly data available to the canonical
kernel, and adjoining them is an allowed pushforward in the bounded-density
lemma.

At stage $s$, write $X_s$ for the current original block and $Y_s$ for
the unprocessed original blocks. The domination invariant remains
\[
 Q_{A_s,X_s,Y_s}\le \exp(b_s)
           (\nu_s\otimes P_{X_s}\otimes P_{Y_s}),\qquad
 b_s\le NK/u+sD_0.
\]
Fresh uniform subset indices, independent Gaussian padding, and retained
past row records are all part of the prefix kernel; none needs an
additional message to future weights. As before, a finite query-and-metric
message costs at most $D_0$ per stage. In particular,
\begin{equation}\label{eq:weighted-average-block}
 \mathbb E_{A_s,Y_s}
 D_{\rm KL}(Q_{X_s\mid A_s,Y_s}\Vert P_{X_s})\le b_s.
\end{equation}
This is the KL chain rule applied to the domination inequality. It is
an averaged bound, so it has no $\log(1/\eta)$ term.

On a successful density test, the orthogonal active projection followed
by independent Gaussian padding has relative entropy at most that of
$X_s$, conditional on $(A_s,Y_s)$. Appending the independent Gaussian
row scalar changes no entropy. All source-sampling randomness used for
these rows is fresh conditional on the prefix and the finite message,
and is independent of the current active matrix. Therefore
Lemma~\ref{lem:weighted-shearer}, followed by the fixed quantizer, gives
\[
 D(Q_{Z_s\mid A_s,Y_s}\Vert\rho_{h,u}^{\otimes k})
 \le \frac{k}{N}
       D(Q_{X_s\mid A_s,Y_s}\Vert P_{X_s})
\]
on the successful branches. On a density-test abort or a later default
branch, generate an unused record from the reference law and ignore it
in the decoder; its entropy contribution is zero. On other approximation failures, keep the already generated row record
even if the decoder subsequently ignores it. No success conditioning
is applied to this record. In particular,
we do not replace an unused row record by a deterministic zero, which
would change its reference-relative entropy.

The first layer has the simpler unconditional bound
\[
 D(Q_{Z_1}\Vert\rho_{1,u}^{\otimes k})
 \le(k/N)D(Q_{W^1}\Vert\mu_{W^1})\le kK/u.
\]
If $N/u>\max\{1,2LD_0/K\}$, then $b_s\le2NK/u$ at every stage.
Since the past records are measurable from $A_s$, conditioning, convexity,
\eqref{eq:weighted-average-block}, and the chain rule now give
\begin{equation}\label{eq:weighted-z-kl}
 \begin{split}
 D(Q_{Z_1,\ldots,Z_{L-1}}\Vert\Pi_u^Z)
 &\le D(Q_{Z_1}\Vert\rho_{1,u}^{\otimes k})
       +\sum_{s=2}^{L-1}\mathbb E
       D(Q_{Z_s\mid A_s,Y_s}\Vert\rho_{h,u}^{\otimes k})\\
 &\le 2LkK/u.
 \end{split}
\end{equation}
For the first term, the conditioning-free first-layer bound is used.
The proof prices the current uniform row subset; it does not assert a
global debit for all conditional active entropies.

\subsection{The remaining finite description}

Only the variance group has size $k$. The pilot, independent correction,
and readout groups together have size $m_b$. Only the dedicated $m_R$ source
positions have nonzero readout coefficients; the other source groups
serve the internal reconstructions. The linear estimator
$\widehat\psi=Fh_{\rm src}$ has identically zero columns on the variance
group, so only its $m_0+m_A\le m_b$ source columns need encoding.
The variance group enters solely through the prescribed average of
its squared activation values. Its coefficient records are the $Z_s$
variables already described.

All other selected rows are stored in the ordinary finite grid. A biased
sample may coincide with a uniform sample or with another biased sample.
It is legitimate simply to store its rounded coefficients again: the
actual code has the corresponding equal values, while a reference code
may have unequal values. The decoder is defined in either case. No
original neuron address, learner-width label, or equality partition is
needed. The reconstruction uses the recorded row functions in their
prescribed source-group positions.

Let $\mathcal J_b$ be the alphabet of all these remaining finite
parameters. The preceding description gives the bound
\begin{equation}\label{eq:weighted-finite-budget}
 \begin{gathered}
 P_b=m_b(d+1)+(L-2)(2q+1)m_b+m_R,\\
 J_b=L\log2+P_b\log(3+4B_{\rm code}/\zeta)
                     +(L-2)q^2\log A_T,\\
 \log|\mathcal J_b|\le J_b.
 \end{gathered}
\end{equation}
The first summand permits abort flags. The coefficient count includes
the biased active rows and Gaussian scalars, the nonzero columns of
$F$, and the readout. The last summand encodes the rounded positive
definite metric matrices. Invalid descriptions use a fixed default
predictor.

Let $J$ denote the resulting finite record. The decoder is a fixed map
of $(J,Z)$ to a predictor. Define the code reference prior
\[
 \Pi_u=\operatorname{Unif}(\mathcal J_b)\otimes\Pi_u^Z.
\]
Importantly, $\Pi_u^Z$ is independent of $J$. Regardless of the actual
dependence between the two code parts,
\[
 D(Q_{J,Z}\Vert\Pi_u)
 =D(Q_Z\Vert\Pi_u^Z)+\log|\mathcal J_b|-H_Q(J\mid Z).
\]
Thus dependence between the adaptively selected finite rows and the
uniform rows costs at most the finite description length, once.

\begin{theorem}[Gaussian-weighted arithmetic code]
\label{thm:weighted-code-kl}
Use the budgets above with $0<u=t^2\le t_0^2$. Suppose
\[
 Q\le \exp(NK/u)\mu_{N,t},\qquad N\ge k,\qquad
 N/u>\max\{1,2LD_0/K\}.
\]
There is a code law $\Gamma_Q$ whose decoded predictor has the same
pointwise internal-disturbance coupling guarantee as
Theorem~\ref{thm:em-compression}, and whose relative entropy obeys
\begin{equation}\label{eq:weighted-kl}
 D(\Gamma_Q\Vert\Pi_u)\le J_b+2LkK/u.
\end{equation}
The code space, decoder, finite alphabet, and quantizer maps are fixed
independently of $N$ and the actual scale $u$; only their Gaussian reference probabilities
depend on $u$.
\end{theorem}
\begin{proof}
The only changes to the approximation construction are the irrelevant
Gaussian padding and sampling the variance group without replacement.
The preceding moment and Hoeffding calculations show that the original
clipping, variance repair, rounding, and disturbance budgets still apply.
Every full virtual layer is clipped and rounded before its variance batch
is selected; hence this batch is a genuine independent uniform subset of
the source layer used in the variance estimate. The prefix enlargement
uses only variables available to its existing kernel and preserves the
bounded-density induction and conditional coupling argument. Thus the
same decoded predictor law has the stated pointwise comparison.
Equations~\eqref{eq:weighted-z-kl} and
\eqref{eq:weighted-finite-budget}, together with the conditional-entropy
identity just above the theorem, prove \eqref{eq:weighted-kl}.
\end{proof}

At fixed ancillary resources, the explicit budgets give
$m_b=\widetilde O(K/e)$ and
$k=\widetilde O(1+t_0^4e^{-2})$.
Consequently the two charges are
\[
 J_b=\widetilde O(LK^2/e^2),\qquad
 2LkK/u=\widetilde O\bigl((LK/u)(1+t_0^4/e^2)\bigr).
\]
The next appendix combines this weighted code at warmer scales with the
ordinary finite code at colder scales, and with the direct posterior
bound at smaller widths.

\section{A common posterior bound over dyadic Gaussian scales}
\label{app:weighted-code}

The Gaussian-weighted code can be tuned to a fixed upper scale.
Tuning the variance sample count to each fixed upper scale keeps the
Gaussian row cost below the finite-parameter cost. A finite sequence
of dyadic scale intervals suffices for this choice. The resulting bound is simultaneous in learner
width and Gaussian scale.

Use the posterior resources in
\eqref{eq:factor-posterior-budgets}, with the product-cylinder cost
$D_*$, so that the retained posterior has density at most
$\exp(NK/t^2)$ relative to its raw Gaussian law. In particular
$\beta=\varepsilon/4$, $\alpha=\varepsilon/12$, and
$e=\alpha r^2/(1000L)<1$.
The rank cap $q$, message charge $D_0$, and the pilot, correction, and
readout counts do not depend on the scale ceiling. Only the variance
count and its consequent scalar mesh need to change.

Put
\begin{equation}\label{eq:weighted-dyadic-scales}
 J_*=\lfloor\log_2(1/e)\rfloor,\qquad
 v_j=2^{-j}\quad(0\le j\le J_*),\qquad B_*=J_*+1.
\end{equation}
Every $u\in(e,1]$ has an index $j$ for which
$v_j/2<u\le v_j$, and every listed $v_j$ lies in $[e,1]$.
For each $j$, use the weighted code with ceiling $t_0=\sqrt{v_j}$.
Write $k_j$ for its variance count, $J_{b,j}$ for its finite-code budget
\eqref{eq:weighted-finite-budget}, $p_j$ for its Gaussian dimension
\eqref{eq:weighted-gaussian-dimension}, and $\Pi_{j,u}$ for its
code reference prior. All these objects are fixed before observing data.
For the cold range $0<u\le e$, let $J_c$ be the ordinary finite-code
budget (EM25) at ceiling $t_0=\sqrt e$.

The information threshold and the common charge are
\begin{equation}\label{eq:weighted-dyadic-budget}
 \begin{split}
 \Delta&=\max\left\{K,\ 2LD_0,\
                  \max_{0\le j\le J_*}\frac{2Kk_j}{v_j}\right\},\\
 \mathcal B&=\max\left\{\Delta+1,\ J_c,\
       \max_{0\le j\le J_*}
       \left[J_{b,j}+\frac{4Lk_jK}{v_j}+1\right]\right\}.
 \end{split}
\end{equation}
Let $M,R_\tau$ be the Gaussian reference-family sizes in
Lemma~\ref{lem:si-grid}, with its parameter $J=\Delta$.
For each dyadic interval also define
\begin{equation}\label{eq:weighted-dyadic-grid}
 R_j=1+\left\lfloor\frac{p_j}{2}\log2\right\rfloor,\qquad
 \tau_{j,\ell}=\sqrt{v_j}\exp(-\ell/p_j)
          \quad(0\le\ell<R_j),
\end{equation}
and put
\[
 R_{\max}=\max_j R_j,\qquad
 \Lambda=\log\frac{3[MR_\tau+B_*R_{\max}+1]}{\delta}.
\]

\begin{theorem}[Weighted-code learning at every mean-field-or-stronger scale]
\label{thm:weighted-all-scales}
Under the model and Boolean-teacher conventions of the paper, the
sample condition
\begin{equation}\label{eq:weighted-sample-budget}
 n\ge\frac4\varepsilon[\mathcal B+\Lambda]
\end{equation}
implies, with probability at least $1-\delta$ over the iid training
sample,
\[
 R(Q_{S,N,t,c})\le\varepsilon
 \qquad\text{simultaneously for every }N\ge C,
                      \quad0<t\le1.
\]
The same statement holds after adding the label-independent prior
restriction $|z_w(x)|>1/2$ at every Boolean input.
\end{theorem}
\begin{proof}
We use one sample-independent reference prior on a tagged predictor
space. Give weight $1/3$ to each of the following parts:
\begin{enumerate}
\item the finite Gaussian family for $NK/t^2\le\Delta$;
\item the uniform law on the ordinary cold finite class;
\item a uniform choice of $j$, then a uniform choice of
$\ell\in\{0,\ldots,R_j-1\}$, followed by the weighted code reference
law $\Pi_{j,\tau_{j,\ell}^2}$.
\end{enumerate}
Use the test-margin indicator loss on the raw Gaussian part and the
ramp loss on both comparison-code parts. These form one bounded loss
on the tagged space. Apply Lemma~\ref{lem:si-risk} once.

Write $u=t^2$ and $b=NK/u$. If $b\le\Delta$, the finite Gaussian
family applies to the actual posterior. Its empirical indicator loss
is zero. Including the mixture weight and confidence term, its charge
is at most $\mathcal B+\Lambda$, so its risk is at most
$2(\mathcal B+\Lambda)/n$.

Suppose $b>\Delta$. Then
$N/u>\max\{1,2LD_0/K\}$, as required by both compression arguments.
The readout-tail estimate retains mass at least $1-\beta$ on the common
robustness radius $r$, and its normalized retained law satisfies the
same density bound with the specified $K$.
If $u\le e$, apply the ordinary code at ceiling $\sqrt e$. Its
relative-entropy charge is at most $J_c$ before mixture weights.

If $u>e$, choose $j$ with $v_j/2<u\le v_j$. From
\eqref{eq:weighted-dyadic-budget},
\[
 N=bu/K>\frac{2k_j u}{v_j}>k_j.
\]
Thus sampling the variance group without replacement is permitted.
Theorem~\ref{thm:weighted-code-kl} gives a law on the $j$th code space
with relative entropy at most
\[
 J_{b,j}+2Lk_jK/u
 \le J_{b,j}+4Lk_jK/v_j.
\]
Choose $\ell=\lfloor p_j\log(\sqrt{v_j}/t)\rfloor$.
Then $0\le\ell<R_j$, $\tau_{j,\ell}\ge t$, and
$p_j\log(\tau_{j,\ell}/t)<1$ (with a non-strict bound also sufficient).
For centered isotropic Gaussians in $p_j$ dimensions,
\[
 \frac{d\mathcal N(0,t^2I)}
      {d\mathcal N(0,\tau_{j,\ell}^2I)}\le\exp(1).
\]
The quantizer and remaining finite-code reference law are fixed within
this dyadic interval. Pushing the domination through the quantizer
therefore changes the relative-entropy charge by at most one. The
mixture and confidence costs are bounded by $\Lambda$.

In either compression branch, the common pointwise coupling gives
empirical ramp loss at most $\alpha$ and
$R(Q_{S,N,t,c})\le\beta+\alpha+R_\psi(\Gamma)$.
Lemma~\ref{lem:si-risk} consequently bounds the desired risk by
\[
 \beta+3\alpha+\frac2n[\mathcal B+\Lambda]
 \le\varepsilon.
\]
All comparisons use the same sample event. Optional global-margin
conditioning retains the same target cylinder, readout-tail bound,
and canonical compression argument, so the proof is unchanged.
\end{proof}

\begin{corollary}[Quartic width and ninth-power depth]
\label{cor:weighted-powers}
Fix $d,c,\varepsilon,\delta$ and put
$F_0=1+\log^2(2+C)$ and $G=\log(2+CL)$.
There is a constant $A_0$ depending only on these fixed quantities such
that
\[
 n\ge A_0 C^4L^9F_0^4G^2
\]
suffices in Theorem~\ref{thm:weighted-all-scales}.
\end{corollary}
\begin{proof}
At fixed input dimension, suppressing only logarithms in the explicit
resources, the finite part of every dyadic code has
\[
 J_{b,j}=\widetilde O(LK^2/e^2).
\]
The finite Gaussian row dimension affects the scale-grid charge only
logarithmically. Since $v_j\in[e,1]$, the explicit variance count obeys
\[
 \frac{k_j}{v_j}
 =\widetilde O\left(\frac1{v_j}+\frac{v_j}{e^2}\right)
 =\widetilde O(e^{-2}).
\]
It follows that
\[
 \Delta=\widetilde O(LK^2/e^2),\qquad
 \mathcal B=\widetilde O(LK^2/e^2).
\]
For the cold code the scale term in (EM26) is only $LK/e$, so it is
also absorbed by this bound. These resource statements keep input
dimension fixed; the exact formulas above retain its full dependence.

For the target-cylinder budgets,
\[
 K=O(CLF_0),\quad e^{-1}=O(CL^3F_0),\quad
 q=O(C^2L^4F_0^2),\quad m_b=O(qG).
\]
The finite parameter count is $O(Lq^2G)$, and its scalar mesh costs
$O(G)$ per parameter. Thus
$J_{b,j}=O(C^4L^9F_0^4G^2)$ uniformly in $j$.
The weighted-row charge and raw cutoff are smaller:
\[
 \max_j Lk_jK/v_j=O(C^3L^8F_0^3G),\qquad
 LD_0=O(C^4L^9F_0^4G).
\]
The ordinary cold code has the same leading finite-parameter order,
and all remaining scale-family charges are logarithmic. Substituting
these estimates in \eqref{eq:weighted-sample-budget} proves the claim.
\end{proof}

\section{A multilevel code for two-step reconstructions}
\label{app:mlmc-code}

This appendix gives another finite class of representatives for the
original robustness support. Its members are finite programs using
arithmetic, tanh, Euclidean projection, and a fixed finite median rule.
They need not themselves be robust networks. All randomness below is
used only to prove the existence of a representative in a fixed finite
class; the representative is deterministic after its coefficients and
finite computation graph have been recorded.

\begin{theorem}[Multilevel finite code]
\label{thm:mlmc-code}
There is a universal constant $A$ with the following property. Let
$d,H\ge1$, $K>0$, and $0<r\le1/6$, and put
\[
 E=K+H+1,\qquad
 T_0=d+\log(2H+2),\qquad
 \Lambda=\log(2+d+E+r^{-1}).
\]
There is a finite class $\mathcal G_{K,r,d,H}$, independent of learner
width, such that
\begin{equation}\label{eq:mlmc-budget}
 \log|\mathcal G_{K,r,d,H}|
 \le A T_0\left[\frac{E^3}{r^3}
                  +\frac{(d+1)E^2}{r^2}\right]\Lambda.
\end{equation}
For every width $N$ and every network $w$ with $\kappa_N(w)\le K$,
there is $g\in\mathcal G_{K,r,d,H}$ whose output is the output of the
original network with input-dependent disturbances of stacked norm at
most $r$. Consequently, if $w\in\mathcal R_{N,c}$ and
$r\le c/\sqrt{1+HE_v(w)}$, then
\[
 \sup_{x\in\{-1,1\}^d}|T(z_w(x))-T(g(x))|\le1/16.
\]
\end{theorem}

We prove the theorem in three steps. First we construct an economical
code for a map containing two activations. We then embed such maps in
a virtual computation of the original network. Finally we choose
spectral cutoffs and error allocations jointly across layers.

\subsection{A two-activation sampling lemma}

Consider the vector-valued map
\begin{equation}\label{eq:mlmc-map}
 F(z,z')=\sum_{i=1}^N w_i\tanh\left(a_i\cdot z+
                \sum_{j=1}^N R_{ij}\tanh(c_j\cdot z')\right),
 \qquad w_i\in\R^{q_o},
\end{equation}
where $a_i\in\R^{q_a}$ and $c_j\in\R^{q_c}$. Assume
\begin{equation}\label{eq:mlmc-map-energy}
 N\sum_i\norm{w_i}_2^2\le k_o,\quad
 \norm{R}_F^2\le k_R,\quad
 \frac1N\sum_i\norm{a_i}_2^2\le k_a,\quad
 \frac1N\sum_j\norm{c_j}_2^2\le k_c.
\end{equation}
Write
\[
 A_i=\sum_j|R_{ij}|,\qquad M_i=\max\{1,A_i\},\qquad
 B=\sum_i\norm{w_i}_2,
\]
\begin{equation}\label{eq:mlmc-work-resources}
 D=\sum_i\norm{w_i}_2M_i,\qquad
 Z=\sum_i\frac{\norm{w_i}_2}{M_i}.
\end{equation}
Cauchy--Schwarz gives
\begin{equation}\label{eq:mlmc-work-energy}
 B\le\sqrt{k_o},\qquad B^2\le ZD,\qquad
 D^2\le2k_o(1+k_R).
\end{equation}
Indeed, $\sum_i A_i^2\le N\norm{R}_F^2$, so
$D\le\sqrt{k_o}(1+\sqrt{k_R})$.

\begin{lemma}[Antithetic map sampling]
\label{lem:mlmc-map-sampling}
Fix $\tau>0$ and a finite collection of evaluation arguments. There
is a random finite program approximating \eqref{eq:mlmc-map} with the
following properties. At each fixed argument, its error is at most
$9\tau$ except on an event of probability at most $e^{-c_1g}$, where
$g$ is the odd number of independent median groups. Its expected number
of stored scalar coefficients is at most
\begin{equation}\label{eq:mlmc-map-work}
 c_2g(1+q_a+q_c+q_o)
                 \frac{k_o(1+k_R)}{\tau^2}.
\end{equation}
After an event of constant probability bounding total coefficient work,
the program can be rounded to a fixed grid with one further error at
most a constant times $\tau$. All parameter ranges and the inverse
grid mesh are polynomial in the energy and dimension upper bounds,
the evaluation-argument norm bounds, $\tau^{-1}$, and the work bound.
No original-width neuron index is needed in the stored program.

The probability assertion also holds conditionally when the evaluation
arguments have been produced by earlier independent random programs.
\end{lemma}

\begin{proof}
We give the sampling, coefficient, and rounding details. Deterministic
feature pruning will first ensure finite coefficient ranges. Removing
roots with $\norm{a_i}_2>B_a$ changes the map by at most
\[
 \sum_{\norm{a_i}_2>B_a}\norm{w_i}_2
       \le\sqrt{k_ok_a}/B_a.
\]
Deleting columns $j$ with $\norm{c_j}_2>B_c$ from $R$ changes it by at
most
\begin{align*}
 \sum_i\norm{w_i}_2\sum_{j:\norm{c_j}_2>B_c}|R_{ij}|
 &\le\sqrt{k_o/N}\,
       \sqrt{\#\{j:\norm{c_j}_2>B_c\}\,k_R}\\
 &\le\sqrt{k_ok_Rk_c}/B_c.
\end{align*}
Choose the two cutoffs to make these errors at most $\tau$ each;
zero numerators require no cutoff. All old energy upper bounds remain
valid after pruning. Recompute $A_i,M_i,B,D,Z$ for the pruned map.
If $B\le\tau$, the zero map already has the required accuracy, so
assume $B>\tau$.

For $A_i>0$, sample $J$ with probability $|R_{iJ}|/A_i$ and put
\[
 X_i=A_i\operatorname{sgn}(R_{iJ})\tanh(c_J\cdot z').
\]
Then $|X_i|\le A_i$ and
$\E X_i=\mu_i:=\sum_jR_{ij}\tanh(c_j\cdot z')$.
For $A_i=0$, use $X_i=0$. Let
\[
 n_{i,0}=\lceil M_i^2\rceil,\qquad n_{i,l}=2^l n_{i,0}.
\]
At level zero define
$\Delta_{i,0}=\tanh(a_i\cdot z+\overline X_{i,n_{i,0}})$.
At level $l\ge1$, split $n_{i,l}$ independent observations into two
equal halves with averages $Y,Y'$, and define
\begin{equation}\label{eq:mlmc-antithetic}
 \Delta_{i,l}=
 \tanh\left(a_i\cdot z+\frac{Y+Y'}2\right)
 -\frac12\tanh(a_i\cdot z+Y)
 -\frac12\tanh(a_i\cdot z+Y').
\end{equation}
All three evaluations reuse the same child activations. The bounded
second derivative of tanh and the fourth moment of a bounded sample
average imply, with universal constants,
\begin{equation}\label{eq:mlmc-level-estimates}
 \E|\Delta_{i,l}|^2\le c_3 2^{-2l},\qquad
 \left|\E\tanh(a_i\cdot z+\overline X_{i,n_{i,l}})
             -\tanh(a_i\cdot z+\mu_i)\right|
       \le c_4 2^{-l}.
\end{equation}
For completeness, Taylor's theorem bounds the absolute value in
\eqref{eq:mlmc-antithetic} by
$\norm{\tanh''}_\infty|Y-Y'|^2/8$.
The fourth moment of a centered average of $m$ independent variables
bounded by $2A_i$ is at most a universal constant times $A_i^4/m^2$.
Since $n_{i,0}\ge A_i^2$, this proves the first estimate. The second
follows from the same Taylor bound and the variance of the average.
The estimates are uniform in both evaluation arguments.

Choose a finite last level $l_*$ and let
\[
 p_i=\frac{\norm{w_i}_2}{M_iZ},\qquad
 \pi_l=\frac{2^{-3l/2}}{\sum_{j=0}^{l_*}2^{-3j/2}}.
\]
Zero-weight roots are omitted. Independently choose a root and level,
and return the vector $Y=w_i\Delta_{i,l}/(p_i\pi_l)$.
The antithetic expectations telescope. Therefore, if $F_*=\E Y$,
\begin{equation}\label{eq:mlmc-root-estimates}
 \norm{F_*-F}_2\le c_4B2^{-l_*},\qquad
 \E\norm{Y}_2^2\le c_5ZD,
 \qquad \E\operatorname{work}(Y)\le c_6D/Z.
\end{equation}
Here scalar work counts child activations, with a fixed constant for
each root. The identities proving the last two estimates are
\[
 \sum_i\frac{\norm{w_i}_2^2}{p_i}=ZD,\quad
 \sum_i p_iM_i^2=D/Z,\quad
 \sum_l2^{-2l}/\pi_l=O(1),\quad
 \sum_l\pi_l2^l=O(1).
\]
Choose $2^{l_*}\ge c_4\sqrt{k_o}/\tau$, or a larger common polynomial
upper bound, making the bias at most $\tau$.

Average $m=\lceil c_7ZD/\tau^2\rceil$ independent copies of $Y$ in
each group. For sufficiently large $c_7$, Chebyshev's inequality makes
the probability of vector error exceeding $\tau$ about $F_*$ at most
$1/64$. The expected work of a group is
\[
 O(D^2/\tau^2+D/Z)=O(D^2/\tau^2),
\]
because $B>\tau$ and $B^2\le ZD$ imply
$D/Z\le D^2/B^2<D^2/\tau^2$.

Use $g$ independent groups and the following fixed median rule: choose
the first group output with at least $(g+1)/2$ group outputs within
distance $4\tau$, and return zero if there is no such group. Except
with probability $e^{-c_1g}$, a strict majority of the groups lie
within $\tau$ of $F_*$. Every such group qualifies, and every selected
group lies within $5\tau$ of $F_*$. If each group output is subsequently
changed by at most $\tau$, the same rule returns a point within
$6\tau$ of $F_*$. This follows because the qualifying ball must
intersect the strict majority of good groups. Thus no Lipschitz property
of the median operation is needed. Combining the two pruning errors,
finite-level bias, and the rounded median gives error at most $9\tau$.

Multiplying scalar work by $1+q_a+q_c+q_o$ bounds the number of recorded
scalar coefficients: record each selected root feature, each child
feature, its signed inner coefficient, and the vector coefficient of
each term in a group average. This proves \eqref{eq:mlmc-map-work}.
For several maps, apply Markov's inequality to their \emph{total
dimension-weighted coefficient work}, rather than separately to their
unweighted child counts. On an event of constant probability, this
total is at most its chosen common deterministic bound $W$. Every
selected root then has $M_i^2\le n_{i,l}\le W$.

The vector coefficient of an averaged term is bounded by
\[
 \norm{\frac{w_i}{mp_i\pi_l}}_2
 \le\frac{\tau^2M_i}{c_7D\pi_l}
 \le\frac{\tau\sqrt W}{c_7\pi_l},
\]
using $D\ge B>\tau$. Inner coefficients have magnitude at most
$A_i\le\sqrt W$. Feature pruning bounded the remaining feature
coefficients, and the finite-level inverse probabilities are polynomial
in the common resources. Thus every recorded coefficient lies in a
fixed polynomially bounded interval.

Suppose all evaluation arguments have norm at most a fixed polynomial
bound. There are at most $W$ stored coefficients and only two nested
activation levels. The 1-Lipschitz property of tanh shows that rounding
all coefficients to a sufficiently fine reciprocal-polynomial mesh
changes every group output by at most $\tau$, uniformly on these
arguments. The preceding median argument applies after this rounding.
The forest topology, its group assignment, and its half-list divisions
have description length $O(W\log(2+W))$. Only numeric feature vectors
are stored, so no original-width neuron label appears. Finally, all
stochastic estimates were uniform at a fixed argument. Conditioning on
arguments produced by independent preceding programs therefore preserves
them, and a union bound applies to any finite collection of arguments.
\end{proof}

\subsection{Realization inside the original robustness ball}

Write $|h|_N=\norm{h}_2/\sqrt N$ and let $k_s=\norm{W^s}_F^2$ for
$2\le s\le H$, with $k_{H+1}=E_v(w)$. Split each internal matrix by
singular value truncation at a threshold $a_s\in(0,1)$:
\[
 W^s=A_s+R_s,\qquad A_s=U_s\Sigma_sV_s^\top,\qquad
 \norm{R_s}_{\rm op}\le a_s,\qquad
 q_s=\operatorname{rank}A_s\le k_s/a_s^2.
\]
For drive coordinates $t_s\in\R^{q_s}$ use the original-space vector
$u_s=\sqrt N U_st_s$. An exact low-rank drive of a tanh state $h$ is
$t_s=\Sigma_sV_s^\top h/\sqrt N$, of norm at most $\sqrt{k_s}$.

The first drive is the one-activation map
$\Sigma_2V_2^\top h^1(x)/\sqrt N$. The next drive uses $h^1(x)$ as its
left boundary. All later drive maps have the two-activation form
\begin{equation}\label{eq:mlmc-drive-map}
 \frac1{\sqrt N}\Sigma_sV_s^\top
 \tanh\left(\sqrt N U_{s-1}t_{s-1}
       +R_{s-1}\tanh(\sqrt N U_{s-2}t_{s-2})\right).
\end{equation}
In Lemma~\ref{lem:mlmc-map-sampling}, the output-energy bound is $k_s$,
the residual-energy bound is $k_{s-1}$, and the root and child feature
second moments are $q_{s-1}$ and $q_{s-2}$. For an early map using the
original first state, its child features are rows of $W^1$, whose
normalized squared norm sum is at most $K/d$. The final readout has
the same form with scalar output and output-energy bound $k_{H+1}$.

Construct and round each map while advancing the computation. After
each drive map, project its output onto the ball of radius
$\sqrt{\max(K,H,1)}$. This does not increase error about its exact
target drive. Let $\widehat t_s$ be the resulting coordinates and
define full virtual states by
\[
 \widehat h^1=h^1,\qquad
 \widehat h^2=\tanh(\sqrt N U_2\widehat t_2+R_2h^1),
\]
\begin{equation}\label{eq:mlmc-virtual-state}
 \widehat h^s=\tanh\left(\sqrt N U_s\widehat t_s+
                  R_s\tanh(\sqrt N U_{s-1}\widehat t_{s-1})\right),
 \qquad s\ge3.
\end{equation}
These full states are used only to verify the disturbance bound; the
stored program advances the drive coordinates through its finite map
codes and evaluates the final readout code.

Suppose the drive-map errors are at most $\delta_s$, and the final
readout-map error is at most $\delta_{H+1}$. Their exact targets are
the low-rank drives of the already rounded preceding virtual states.
For $s\ge3$, comparison with the original recurrence gives
\[
 |\widehat h^s-\tanh(W^s\widehat h^{s-1})|_N
 \le\delta_s+a_sa_{s-1}.
\]
Indeed, the low-rank preactivation discrepancy has normalized norm
at most $\delta_s$, while the discarded older residual passes through
$R_{s-1}$ and $R_s$, with a 1-Lipschitz tanh between. Its left boundary
is a tanh state of normalized norm at most one. At stage two the bound
is just $\delta_2$; the first-state disturbance is zero. Consequently
the original-network disturbances realizing the virtual output satisfy
\begin{equation}\label{eq:mlmc-disturbance}
 \norm\eta_*
 \le\left(\sum_{s=2}^{H+1}\delta_s^2\right)^{1/2}
   +\left(\sum_{s=3}^H a_s^2a_{s-1}^2\right)^{1/2}.
\end{equation}
For $H=1$, only the one-activation readout map is needed; for $H=2$,
the exact first-state boundary is retained throughout.

Every stage uses fresh map samples independent of earlier map programs.
Lemma~\ref{lem:mlmc-map-sampling} therefore applies conditionally at
the already rounded prefix. Choose $g=O(T_0)$ so a union bound over
all $2^d$ inputs and all maps fails with probability at most $1/100$.
The coefficient-work event can be chosen to fail with probability at
most $1/100$, so the two events intersect. All rounding is already
included in the map errors in \eqref{eq:mlmc-disturbance}.

\subsection{Adaptive cutoffs and a common finite class}

Put $e_s=1+k_s$ for the internal blocks and the readout. Their sum is
at most $E=K+H+1$. For a sufficiently small universal constant $c_0$,
choose
\begin{equation}\label{eq:mlmc-cutoffs}
 a_s^2=\frac{c_0r\sqrt{e_s}}{\sqrt E},\qquad 2\le s\le H.
\end{equation}
Then $a_s<1$ and
\[
 \sum_{s=3}^H a_s^2a_{s-1}^2
 =\frac{c_0^2r^2}{E}\sum_{s=3}^H\sqrt{e_se_{s-1}}
 \le c_0^2r^2.
\]
With $\alpha=\sqrt E/(c_0r)$, the retained ranks satisfy
$q_s\le\alpha\sqrt{e_s}$.

Ignoring fixed constants, define each map's coefficient resource as
\begin{equation}\label{eq:mlmc-profile-resource}
 B_s=(1+q_{s-2}+q_{s-1}+q_s)k_s(1+k_{s-1}).
\end{equation}
At the first two maps, replace the missing input dimension by $d+1$
and omit a nonexistent residual factor. At the readout, replace the
output dimension by one and the output energy by $k_{H+1}$. These
conventions give exactly the resources in
Lemma~\ref{lem:mlmc-map-sampling}. Apart from the first two original-
input contributions, they are bounded by a universal constant times
\[
 \alpha e_se_{s-1}
    (\sqrt{e_{s-2}}+\sqrt{e_{s-1}}+\sqrt{e_s}).
\]
At boundary indices in this estimate, set a missing $e_j$ equal to one;
the finitely many auxiliary indices change the following bounds only
by universal constants.
For each of the three indicated shifts $j$,
\[
 \sum_s e_s^{1/2}e_{s-1}^{1/2}e_j^{1/4}
 \le(\max_j e_j)^{1/4}\sum_s\sqrt{e_se_{s-1}}
 \le c_8' E^{5/4}.
\]
Using $\sqrt{u+v+w}\le\sqrt u+\sqrt v+\sqrt w$, and treating the
two original-input terms separately, gives
\begin{equation}\label{eq:mlmc-profile-sum}
 \left(\sum_s\sqrt{\max(1,B_s)}\right)^2
 \le c_8\left[\frac{E^3}{r}+(d+1)E^2\right].
\end{equation}
Inserting the ones increases the unsquared sum by at most $H+1$,
so its additional squared contribution is $O(H^2)$, already absorbed
by the right side. Each original-input term contributes at most a constant
times $\sqrt{d+1}\,E$ before squaring.

Let $S_0$ be the square root of the right side of
\eqref{eq:mlmc-profile-sum}, enlarged by a fixed constant. Set
\[
 \delta_s^2=\frac{r^2}{64S_0}\sqrt{\max(1,B_s)},
\]
and round each $\delta_s$ downward to a power-of-two multiple of a fixed
reciprocal-polynomial base, losing at most a factor two. Thus
$\sum_s\delta_s^2\le r^2/64$, every nonzero tolerance has a common
reciprocal-polynomial lower bound, and
\begin{equation}\label{eq:mlmc-total-work}
 \sum_s B_s\delta_s^{-2}
 \le c_9\left[\frac{E^3}{r^3}
                    +\frac{(d+1)E^2}{r^2}\right].
\end{equation}
Apply the sampling lemma with $\tau_s$ a sufficiently small fixed
multiple of $\delta_s$. Equations \eqref{eq:mlmc-disturbance} and
\eqref{eq:mlmc-cutoffs} put the complete virtual computation inside
the radius-$r$ disturbance ball.

To make the code common to all widths and all retained networks, record
the integer rank profile, the dyadic tolerance choices, and the finite
map forests. Every rank is at most $E/(c_0r)$. The bounds
\eqref{eq:mlmc-total-work} and $g=O(T_0)$ give a common polynomial bound
$W$ on total coefficient work on the chosen Markov event. Feature
cutoffs, coefficient magnitudes, inverse grid mesh, and finite last
levels are all polynomial in $2+d+E+r^{-1}$, with universal degrees:
this follows from the common rank bound, the common lower tolerance
bound, and the coefficient estimates in the sampling lemma. Fix those
ranges and that grid before defining the class of all recorded programs.
The possible topology, ranks, and tolerance choices cost
$O(W\Lambda+H\Lambda)$ nats; the coefficient grid costs $O(W\Lambda)$.
Since $H\le E$ and $r\le1/6$, the former additional term is absorbed.
This proves \eqref{eq:mlmc-budget}. The class contains no learner-width labels,
so the same class works for every $N$. This completes the proof of
Theorem~\ref{thm:mlmc-code}.

\subsection{The original posterior at the stronger scales}
\label{sec:mlmc-posterior}

The new code changes neither the prior nor the conditioning event.
The following statement records the exact scale hypotheses used with it.
Use the product-cylinder cost
$D_*=D_*^{\rm prod}$ from \eqref{eq:product-cylinder-cost}, and
write
\[
 P_N=N(d+1)+(L-2)N^2+N.
\]
For the chosen energy cutoff $K$, retain the radius
\[
 r=\min\{1/6,c/\sqrt{1+HK}\}.
\]
Let $J_{\rm ML}$ be the upper bound in \eqref{eq:mlmc-budget}, and let
$J$ be the earlier dense-code budget in \eqref{eq:budget}, with this same
$K$ and $r$. Put
\[
 J_{\rm best}=\min\{J,J_{\rm ML}\}.
\]
This minimum is legitimate without a union penalty: choose, from the
known resource bounds, whichever of the two common finite classes has
the smaller stated cardinality bound.

\begin{theorem}[A mixed width--depth bound at the original stronger scales]
\label{thm:mlmc-posterior}
Fix $0<\varepsilon,\delta<1$. Suppose the target has the Boolean
implementation specified in the paper. The sample condition
\[
 n\ge\frac2\varepsilon
       \left[J_{\rm best}+\log(1/\delta)\right]
\]
implies, with sample probability at least $1-\delta$,
$R(Q_{S,N,t,c})\le\varepsilon$ simultaneously at every indicated width
and scale in either of the following regimes:
\begin{enumerate}
\item Set $K=8D_*$. Allow every $N\ge C$ and every positive scale with
\[
 t^2\le\min\left\{1,
 \frac{ND_*}{(P_N/2)\log2+\log(2/\varepsilon)}\right\}.
\]
\item Set
\[
 K=8D_*+4(d+L)+4\log(2/\varepsilon).
\]
Allow every $N\ge C$ and $0<t\le N^{-1/2}$, in particular every
power-law scale with $\gamma\ge2$.
\end{enumerate}
In the first regime, at each fixed width the same sample event gives
$\limsup_{t\downarrow0}R(Q_{S,N,t,c})\le\varepsilon$.
At fixed $d,c,\varepsilon,\delta$, a sufficient sample count is
\[
 \widetilde O\!\left(
      \min\{C^4L^7,C^{9/2}L^6\}\right).
\]
\end{theorem}

\begin{proof}
Every robust network with $\kappa_N(w)\le K$ has calibrated radius at
least $r$. Theorem~\ref{thm:mlmc-code}, or the earlier dense code,
therefore supplies a representative with clipped-score error at most
$1/16$ at every Boolean input. The representative need not be a robust
network; only its membership in a common finite class is used next.

Use the ramp loss that is one at signed clipped scores at most $1/3$
and zero at scores at least $5/12$. Every representative of a fitting
network has zero empirical ramp loss, because its signed clipped
training score is at least $7/16$. At every test-margin failure of the
original network, its representative has signed clipped score at most
$5/16$, so the ramp loss is one. A union bound over the chosen finite
class shows that, on a single sample event of probability at least
$1-\delta$, every retained fitting network has test-margin error at
most $\varepsilon/2$. This sample event is independent of learner width
and Gaussian scale.

The product-cylinder lower bound and the Gaussian exponential moment give
the same tail as \eqref{eq:tail}, with the current valid cylinder cost:
\[
 Q_{S,N,t,c}\{\kappa_N>K\}
 \le\exp\!\left[-\frac{NK}{4t^2}
       +\frac{P_N}{2}\log2+\frac{ND_*}{t^2}\right].
\]
In the first regime the displayed scale inequality makes this at most
$\varepsilon/2$. In the second regime use
$t^2P_N/N\le d+L$ and the stated value of $K$, exactly as in the proof
of Theorem~\ref{thm:main}. Adding the discarded mass proves the risk
claim. At fixed $N$, the first scale inequality holds for all sufficiently
small positive $t$, proving the stated strong-prior conclusion without
identifying a separately defined zero-temperature posterior.

Finally, the product-cylinder calculation gives
$K=\widetilde O(CL)$ and $r^{-2}=\widetilde O(CL^2)$.
The budget \eqref{eq:mlmc-budget}, with $K+H+1=\widetilde O(CL)$,
is therefore $\widetilde O(C^{9/2}L^6)$ at the fixed resources stated
in the theorem. Taking the better of this and the existing
$\widetilde O(C^4L^7)$ code gives the final display.
\end{proof}

\begin{corollary}[Explicit logarithms in the mixed-power budget]
\label{cor:mlmc-powers}
Fix $d,c,\varepsilon,\delta$, and put
$F_0=1+\log^2(2+C)$ and $G=\log(2+CL)$.
In either regime of Theorem~\ref{thm:mlmc-posterior}, a sufficient
sample count is
\[
 O\!\left(G\min\left\{
 C^4L^7F_0^4\log^2(2+L),\;
 C^{9/2}L^6F_0^{9/2}\log(2+L)
 \right\}\right).
\]
In particular, at fixed teacher width the new branch is
$O(L^6\log^2(2+L))$.
\end{corollary}

\begin{proof}
The product-cylinder cost satisfies $D_*^{\rm prod}=O(CLF_0)$.
For either stated energy cutoff,
\[
 E=O(CLF_0),\qquad r^{-2}=O(CL^2F_0),\qquad
 T_0=O(\log(2+L)),\qquad \Lambda=O(G).
\]
Consequently the leading term in \eqref{eq:mlmc-budget} is
\[
 O\!\left(C^{9/2}L^6F_0^{9/2}\log(2+L)G\right).
\]
The term $(d+1)E^2r^{-2}$ has order $O(C^3L^4F_0^3)$ and is smaller
at the fixed ancillary resources. The dense-code calculation with the
same product-cylinder substitution gives
$O(C^4L^7F_0^4\log^2(2+L)G)$, as recorded after
Corollary~\ref{cor:all-scales-product-cylinder}. Taking the better code
and applying Theorem~\ref{thm:mlmc-posterior} proves the result.
\end{proof}

The second branch improves the stated depth dependence at fixed teacher
width, and is smaller at the level of the displayed powers when
$C<L^2$. It does not establish a better uniform quartic-width exponent.
No bounded spectral product, column coherence, additional output-margin
condition, or modified robustness support is used.

\section{A width-independent positive cold-scale range}
\label{app:spectral-cold}

An operator-norm truncation separates a small Gaussian background from
a bounded-energy network. The background need not have small Frobenius
norm: its operator norm controls its effect at every bounded virtual
hidden state. The excess above that operator-norm threshold has Gaussian
concentration independent of learner width. This gives a positive range
of scales on which the stronger energy-based sample bounds hold at
\emph{all} learner widths, rather than requiring $t^2\le1/N$.

\subsection{Gaussian spectral excess and its posterior tail}

For a square matrix $W$ and $a\ge0$, let $\mathcal B_a$ be the closed
operator-norm ball of radius $a$. If the singular values of $W$ are
$\sigma_i$, its Frobenius projection onto $\mathcal B_a$ clamps each
singular value to $\min\{\sigma_i,a\}$. The residual therefore has
singular values $(\sigma_i-a)_+$. Fix the absolute threshold constant
\[
 a_0=8.
\]
For a network $w$, leave the first layer and readout unchanged and
replace each internal matrix by its residual above $a_0t$:
\[
 V^1=W^1,\qquad v_V=v,\qquad
 V^s=W^s-B^s,\quad
 B^s=\operatorname{Proj}_{\mathcal B_{a_0t}}(W^s)
              \quad(2\le s\le H).
\]
In particular $\|B^s\|_{\rm op}\le a_0t$.
Define the spectral-excess radius
\begin{equation}\label{eq:spectral-radius}
 \mathcal E_t(w)^2=\kappa_N(V)
 =\frac dN\|W^1\|_F^2+N\|v\|_2^2
       +\sum_{s=2}^H
            \operatorname{dist}_F(W^s,\mathcal B_{a_0t})^2.
\end{equation}

\begin{lemma}[A width-independent spectral-excess estimate]
\label{lem:spectral-excess}
Under the raw Gaussian law $\mu_{N,t}$, with $u=t^2$,
\[
 \mathbb E\mathcal E_t^2\le u(d+2H),
\]
and for every $s\ge0$,
\begin{equation}\label{eq:spectral-concentration}
 \mu_{N,t}\{\mathcal E_t\ge\mathbb E\mathcal E_t+s\}
       \le\exp\left(-\frac{Ns^2}{2u}\right).
\end{equation}
\end{lemma}
\begin{proof}
A pair of $1/4$-nets of the unit sphere, each of cardinality at most
$9^N$, gives for an internal Gaussian matrix the elementary bound
\[
 \Pr\{\|W\|_{\rm op}>x\}
 \le2\exp\left(2N\log9-\frac{Nx^2}{8u}\right).
\]
Indeed, the operator norm is at most twice the largest bilinear form
on the two nets, and each such form is Gaussian of variance $u/N$.
Writing $R=\operatorname{dist}_F(W,\mathcal B_{8t})$, we have
$R^2\le N(\|W\|_{\rm op}-8t)_+^2$. Hence
\begin{align*}
 \mathbb E R^2
 &\le 4N\int_0^\infty s
       \exp\left(2N\log9-
                   \frac{N(8t+s)^2}{8u}\right)ds\\
 &\le 16u\exp[-N(8-2\log9)]\le u.
\end{align*}
The last inequality follows from $N\ge1$ and
$16\exp[-(8-2\log9)]<1$.
The expected first-layer contribution to \eqref{eq:spectral-radius}
is $u(d+1)$ and the readout contribution is $u$.
There are $H-1$ internal matrices. Thus the expectation is at most
$u(d+H+1)\le u(d+2H)$.

In the weighted coordinates
\[
 y=\bigl(\sqrt{d/N}\,W^1,\ W^2,\ldots,W^H,\ \sqrt N\,v\bigr),
\]
the raw law has independent coordinates of common variance $u/N$.
The quantity $\mathcal E_t$ is Euclidean distance to the closed convex
product set whose first and readout coordinates are zero and whose
internal blocks belong to $\mathcal B_{8t}$. Distance to a nonempty
closed set is 1-Lipschitz. Thus $\mathcal E_t$, as a function of a
standard Gaussian vector, is $t/\sqrt N$-Lipschitz.

For completeness, the Gaussian Lipschitz concentration estimate used
here follows from the log-Sobolev inequality and Herbst's calculation.
For a positive smooth function $g$, Ornstein--Uhlenbeck entropy
dissipation and score commutation give
\[
 \operatorname{Ent}_{\gamma}(g)
 \le\frac12\mathbb E_\gamma\frac{\|\nabla g\|^2}{g}.
\]
For $g=e^{\lambda F}$ and a function $F$ of Lipschitz constant $\sigma$,
this implies
$\log\mathbb E e^{\lambda(F-\mathbb EF)}\le\lambda^2\sigma^2/2$:
if $\chi(\lambda)=\log\mathbb Ee^{\lambda F}$, the entropy inequality
says $\lambda\chi'-\chi\le\lambda^2\sigma^2/2$, and integration
from zero proves the assertion. Markov's inequality and optimization in
$\lambda$ give the upper tail $e^{-s^2/(2\sigma^2)}$.
Smooth approximation gives the same result for Lipschitz distance
functions. Substituting $\sigma=t/\sqrt N$ proves
\eqref{eq:spectral-concentration}.
\end{proof}

Let $D_*=D_*^{\rm prod}$ be the product target-cylinder cost. The actual
posterior, for every correctly labeled training sample, obeys
\[
 Q_{S,N,t,c}\le\exp(ND_*/u)\mu_{N,t}.
\]
Fix $0<\varepsilon<1$ and set
\begin{equation}\label{eq:spectral-budget}
 \begin{split}
 \beta&=\varepsilon/2,\\
 K_{\rm sp}&=2(d+2H)+4[D_*+\log(1/\beta)],\\
 r_{\rm sp}&=\min\{1/6,c/\sqrt{1+HK_{\rm sp}}\},
 \qquad \rho_{\rm sp}=r_{\rm sp}/2.
 \end{split}
\end{equation}
Then, uniformly over all $N\ge C$ and $0<u\le1$,
\begin{equation}\label{eq:spectral-posterior-tail}
 Q_{S,N,t,c}\{\kappa_N(V)>K_{\rm sp}\}\le\beta.
\end{equation}
To see this, put $b_0=d+2H$ and $A=D_*+\log(1/\beta)$.
The inequality
$\sqrt{2b_0+4A}\ge\sqrt{b_0}+\sqrt{2A}$ and the preceding lemma imply
\[
 \mu_{N,t}\{\mathcal E_t^2>K_{\rm sp}\}
 \le \exp(-NA/u).
\]
Multiply by the posterior density bound. Since $N/u\ge1$, the result
is at most $\exp[-N\log(1/\beta)/u]\le\beta$.

\subsection{Combining spectral truncation with an energy code}

On the retained event, the original readout satisfies
$E_v=N\|v\|^2\le\kappa_N(V)\le K_{\rm sp}$.
Every original network in the posterior support is therefore robust on
the common radius $r_{\rm sp}$. No robustness of the residual network
$V$ is required.

Apply either the dense sampling code or the multilevel code to $V$,
with energy cutoff $K_{\rm sp}$ and disturbance budget
$\rho_{\rm sp}=r_{\rm sp}/2$. Both constructions have full virtual
hidden states $\widehat h^s(x)\in[-1,1]^N$ and reproduce their decoded
output by injections $\eta$ into $V$, with
$\|\eta\|_*\le\rho_{\rm sp}$.
For the dense construction the proved bound is even
$\rho_{\rm sp}/2$; the weaker common bound suffices here.
At these \emph{same} virtual states define injections into the original
network by
\[
 \xi^1=\eta^1,\qquad
 \xi^s=\eta^s+
   \tanh(V^s\widehat h^{s-1})-
   \tanh(W^s\widehat h^{s-1})\quad(2\le s\le H),
 \qquad \xi^{\rm out}=\eta^{\rm out}.
\]
The readout is unchanged, and these injections reproduce exactly the
same states and decoded output. Since tanh is 1-Lipschitz,
\[
 \frac1{\sqrt N}\left\|
 \tanh(V^s\widehat h^{s-1})-
 \tanh(W^s\widehat h^{s-1})\right\|_2
 \le\|B^s\|_{\rm op}\le8t.
\]
Therefore
\[
 \|\xi\|_*\le\rho_{\rm sp}+8t\sqrt{H-1}.
\]
The following positive scale ceiling is sufficient:
\begin{equation}\label{eq:spectral-cutoff}
 u_c=\frac{r_{\rm sp}^2}{256H}.
\end{equation}
For $t^2\le u_c$, the right side is at most $r_{\rm sp}$.
The original support condition consequently gives clipped-score error
at most $1/16$ for the decoded representative at every Boolean input.
The representative and the residual network $V$ need not satisfy that
support condition themselves.

Let $J_{\rm den}(K,\rho)$ denote the dense-code budget in
\eqref{eq:budget}, recalculated with energy $K$ and radius $\rho$.
Let $J_{\rm ML}(K,\rho)$ denote the bound in
\eqref{eq:mlmc-budget} at those same resources. Put
\begin{equation}\label{eq:spectral-code-budget}
 J_{\rm sp}=\min\{J_{\rm den}(K_{\rm sp},\rho_{\rm sp}),
                    J_{\rm ML}(K_{\rm sp},\rho_{\rm sp})\}.
\end{equation}
Choose whichever finite comparison class has the smaller known upper
bound; this choice is made before the sample is drawn and needs no
extra union bound.

\begin{theorem}[Positive cold scales uniformly over learner width]
\label{thm:spectral-cold}
For the Boolean teacher, Gaussian tanh learner, and unchanged robustness
support of the paper, suppose
\[
 n\ge\frac2\varepsilon\,[J_{\rm sp}+\log(1/\delta)].
\]
Then, with sample probability at least $1-\delta$,
\[
 R(Q_{S,N,t,c})\le\varepsilon
\]
simultaneously for every $N\ge C$ and every $0<t^2\le u_c$, with
$K_{\rm sp},r_{\rm sp},u_c,J_{\rm sp}$ as above.
The ceiling $u_c$ is positive and independent of learner width.
The same conclusion holds with optional label-independent global-margin
conditioning.
\end{theorem}
\begin{proof}
For every retained original network that fits the training margins,
the preceding construction supplies a representative in the common
chosen finite class. Its clipped score differs by at most $1/16$ at
every input. Its empirical ramp loss is therefore zero, and its ramp
loss is one at every test-margin failure of the original network.
For any fixed representative with mean ramp loss greater than
$\varepsilon/2$, the probability of zero empirical ramp loss is at
most $\exp(-n\varepsilon/2)$. A union bound over the finite class shows
that every fitting retained original network has test-margin error
at most $\varepsilon/2$ on one sample event of probability at least
$1-\delta$. This event is independent of $N,t$.
Adding the discarded posterior mass in
\eqref{eq:spectral-posterior-tail}, at most $\beta=\varepsilon/2$,
proves the claim. The optional initial global-margin condition keeps
the same target-cylinder denominator and the same support, so all
steps remain valid.
\end{proof}

At fixed $d,c,\varepsilon,\delta$, the product-cylinder calculation gives
\[
 K_{\rm sp}=O(CLF_0),\qquad
 r_{\rm sp}^{-2}=O(CL^2F_0),\qquad
 u_c^{-1}=O(CL^3F_0),
 \qquad F_0=1+\log^2(2+C).
\]
Thus the same sufficient sample order as in the stronger-scale energy
result holds throughout this width-independent cold interval:
\[
 \widetilde O\!\left(\min\{C^4L^7,C^{9/2}L^6\}\right).
\]
Its log-retaining form, with $G=\log(2+CL)$, is
\[
 O\!\left(G\min\left\{
 C^4L^7F_0^4\log^2(2+L),\quad
 C^{9/2}L^6F_0^{9/2}\log(2+L)
 \right\}\right).
\]
The finite-scale statement is stronger than merely taking a limit at
each fixed width: it provides a positive scale threshold common to
all widths above target complexity.

\section{Interpolating the width-uniform cold bound}
\label{app:spectral-interpolation}

The spectral residual argument gives a finite-scale theorem with a
threshold that does not depend on learner width. It can also sharpen
the dependence on the exponent $\gamma$ before the regime $\gamma\ge2$.
We record the interpolation explicitly, including its common sample
event.

Let $0<u_c\le1$ be a width-independent cold threshold, and let
$\mathcal G_c$ be its finite comparison class, with
$\log|\mathcal G_c|\le J_c$. Suppose the cold construction has the
following consequences for every $N\ge C$ and $0<t^2\le u_c$:
there is posterior mass at least $1-\beta$ whose members have
representatives in $\mathcal G_c$ with zero empirical ramp loss and
population ramp loss at least their test-margin error. In particular,
a law $\Gamma_c$ on $\mathcal G_c$ can be chosen with
\begin{equation}\label{eq:spectral-cold-interface}
 \widehat R_{\psi,S}(\Gamma_c)=0,\qquad
 R(Q_{S,N,t,c})\le\beta+R_\psi(\Gamma_c).
\end{equation}
The original posterior still satisfies
$Q_{S,N,t,c}\le\exp(NK/t^2)\mu_{N,t}$, where $K\ge1$ is any
valid target-cost bound, for example $D_*+1$.
Theorem~\ref{thm:spectral-cold} supplies this interface with the
threshold in \eqref{eq:spectral-cutoff}, the code budget in
\eqref{eq:spectral-code-budget}, and $\beta=\varepsilon/2$. The raw
target cost $K$ here is distinct from its residual-energy cutoff.

\begin{corollary}[Exponent interpolation from a uniform cold threshold]
\label{cor:spectral-interpolation}
Assume \eqref{eq:spectral-cold-interface}, and fix $\gamma>1$.
Put
\[
 B_\gamma=K u_c^{-\gamma/(\gamma-1)}.
\]
Let $M_\gamma,R_{\tau,\gamma}$ be the finite Gaussian reference-family
sizes of Lemma~\ref{lem:si-grid}, with its parameter $J=B_\gamma$, and
set
\[
 \Lambda_\gamma=1+
     \log\frac{2M_\gamma R_{\tau,\gamma}}{\delta}.
\]
With sample probability at least $1-\delta$, simultaneously for every
$N\ge C$ and $0<t^2\le N^{1-\gamma}$,
\begin{equation}\label{eq:spectral-interpolation-risk}
 R(Q_{S,N,t,c})\le\beta+
             \frac2n[J_c+B_\gamma+\Lambda_\gamma].
\end{equation}
Consequently, if $\beta\le\varepsilon/2$, it suffices to take
\begin{equation}\label{eq:spectral-interpolation-samples}
 n\ge\frac4\varepsilon
       [J_c+B_\gamma+\Lambda_\gamma].
\end{equation}
One may replace the sum $J_c+B_\gamma$ by their maximum after
adjusting the harmless finite-family confidence charge.
\end{corollary}
\begin{proof}
Use one reference prior with two equally weighted parts: the uniform
law on $\mathcal G_c$, and the finite Gaussian family from
Lemma~\ref{lem:si-grid}. Use the ramp loss on the first component and
the test-margin indicator on the Gaussian components. Apply the single
relative-entropy risk event of Lemma~\ref{lem:si-risk}.

Write $u=t^2$. If $u\le u_c$, the law in
\eqref{eq:spectral-cold-interface} has zero empirical loss and relative
entropy at most $J_c+\log2$ under this reference prior. Its population
loss is at most $2[J_c+\log(2/\delta)]/n$.

If $u>u_c$ and $u\le N^{1-\gamma}$, then
\[
 N\le u^{-1/(\gamma-1)},\qquad
 NK/u\le K u^{-\gamma/(\gamma-1)}<B_\gamma.
\]
The Gaussian reference family therefore applies to the actual posterior.
Its empirical margin-indicator loss is zero, and its relative-entropy
charge is at most
$B_\gamma+1+\log(2M_\gamma R_{\tau,\gamma})$.
The resulting risk bound is at most
$2[B_\gamma+\Lambda_\gamma]/n$.
These two cases occur on the same sample event and prove
\eqref{eq:spectral-interpolation-risk}. The sample condition follows
by allocating at most $\varepsilon/2$ to the posterior tail and at
most $\varepsilon/2$ to the statistical term.
\end{proof}

For the Boolean target-cylinder budgets and the spectral residual
construction, at fixed $d,c,\varepsilon,\delta$,
\[
 K=\widetilde O(CL),\qquad
 u_c^{-1}=\widetilde O(CL^3),\qquad
 J_c=\widetilde O(C^4L^7).
\]
Thus for every fixed $\gamma>1$, the interpolated sufficient order is
\begin{equation}\label{eq:spectral-interpolation-powers}
 \widetilde O\left(
 C^4L^7+
 C^{\,2+1/(\gamma-1)}L^{\,4+3/(\gamma-1)}\right).
\end{equation}
Taking the better of this and the all-scale weighted-code theorem gives
\[
 \widetilde O\!\left(\min\left\{
 C^4L^9,\;
 C^4L^7+C^{\,2+1/(\gamma-1)}L^{\,4+3/(\gamma-1)}
 \right\}\right).
\]
The choice of which sufficient bound to use can be made from the resource
budgets before the sample is drawn. No sample-dependent choice of a
reference prior is required.

In particular, $\gamma\ge7/4$ permits
$\widetilde O(C^4L^8)$, and $\gamma\ge2$ permits
$\widetilde O(C^4L^7)$. These statements can use constants uniform in
the indicated exponent ranges. At the endpoint $\gamma=7/4$, the
more precise two-term order is
$\widetilde O(C^4L^7+C^{10/3}L^8)$.
The separate multilevel bound
$\widetilde O(\min\{C^4L^7,C^{9/2}L^6\})$ remains available in its
stated stronger-scale regimes.

For explicit logarithms, write $F_0=1+\log^2(2+C)$ and
$G=\log(2+CL)$. The target and cold-threshold estimates give
\[
 B_\gamma=
 O\!\left(C^{\,2+1/(\gamma-1)}
          L^{\,4+3/(\gamma-1)}F_0^{\,2+1/(\gamma-1)}\right),
\]
with fixed $\gamma>1$ and fixed ancillary resources.
For $\gamma\ge7/4$, a common log-retaining upper bound is
\[
 O\!\left(C^4L^8F_0^4\log^2(2+L)G\right).
\]
The finite-family charge $\Lambda_\gamma$ is logarithmic in the
resource budgets; it is uniformly $O(G)$ when $\gamma\ge7/4$ at the
fixed confidence parameter. The all-scale theorem has its own
$F_0^4G^2$ factor and remains preferable when the displayed interpolation
is larger.


\begin{thebibliography}{9}
\bibitem{sequential}
ChatGPT-6 Astra, Claude Fable 5 and Opus 5, Kaarel H\"anni, and Dmitry Vaintrob.
\emph{Sharper depth dependence for Bayesian networks stable under internal noise}.
Research manuscript, 29 October 2026.
\href{https://github.com/kaarelh/AI-works/tree/main/neural-network-generalization/works-in-progress/sequential-robust-quantization}{Source and build instructions}.
\bibitem{quartic}
ChatGPT-6 Astra, Claude Fable 5 and Opus 5, Kaarel H\"anni, and Dmitry Vaintrob.
\emph{A quartic-width learning bound for Bayesian networks stable under internal noise}.
Research manuscript, 21 October 2026.
\href{https://github.com/kaarelh/AI-works/tree/main/neural-network-generalization/works-in-progress/improved-polynomials/original-support}{Source and build instructions}.
\bibitem{robust-paper}
ChatGPT-6 Astra, Claude Fable 5 and Opus 5, Kaarel H\"anni, and Dmitry Vaintrob.
\emph{A polynomial-depth learning bound for NNBayes on robust networks}.
Research manuscript, 18 October 2026.
\href{https://github.com/kaarelh/AI-works/tree/main/neural-network-generalization/main-results/robust-polynomial-depth}{Source and build instructions}.
\bibitem{all-scales}
ChatGPT-6 Astra, Claude Fable 5 and Opus 5, Kaarel H\"anni, and Dmitry Vaintrob.
\emph{A sharper polynomial learning bound for NNBayes on robust networks}.
Research manuscript, 21 October 2026.
\href{https://github.com/kaarelh/AI-works/tree/main/neural-network-generalization/works-in-progress/improved-polynomials/original-support-all-scales}{Source and build instructions}.
\bibitem{tarres-yao}
Pierre Tarr\`es and Yuan Yao.
\emph{Online learning as stochastic approximation of regularization paths:
optimality and almost-sure convergence}.
Appendix A, Corollary A.2, equation (A-1).
\url{https://arxiv.org/abs/1103.5538}.
\end{thebibliography}
\end{document}
